%Written by Rishal S P 
%Almost the entire document's text was edited slightly and rewritten based on suggestions by Notion AI to improve clarity.
%This was done on 18th Oct 2024 and 21st Oct 2024. 

\documentclass[acmsmall, screen, nonacm, anonymous]{acmart}
%review,
%\documentclass[acmsmall,nonacm,anonymous]{acmart}

\usepackage[capitalise]{cleveref}
\usepackage{graphicx}
\let\Bbbk\relax
\usepackage{amsmath, amssymb}
\usepackage{amsthm}

%\usepackage[hidelinks]{hyperref}
\usepackage{algorithm}
\usepackage{algpseudocodex} 
%both packages for making algorithm figs
\usepackage{caption}
\usepackage{todonotes} %[disable]
%\usepackage[appendix=inline]{apxproof}
\usepackage{apxproof}
%\renewcommand{\appendixprelim}{\clearpage \section*{Appendix}} 
\usepackage{tikz}
 \usepackage{tikzit}
 \input{figures/styles.tikzstyles}


\bibliographystyle{alpha}% the mandatory bibstyle

\title{Krama: Participant-Centric Consensus for Blockchains}

%Listed Alphabetically on last name
%\author{Rahul Lenkala}{Sarva Labs}{rahul@sarva.ai}{}{}

%\author{Srinidhi Nagendra}{Max Planck Institute for Software Systems, Germany}{hello@srinidhin.com}{}{}

%\author{S~P~Rishal}{Sarva Labs}{rishal@sarva.ai}{}{}

%\Copyright{Rahul, Srinidhi, and S~P} %TODO mandatory, please use full first names. LIPIcs license is "CC-BY";  http://creativecommons.org/licenses/by/3.0/

%\keywords{Blockchain, Byzantine-fault tolerance, State machine replication, Consensus, Participant-centric consensus, Context }
%
%\category{} %optional, e.g. invited paper
%
%\relatedversion{} %optional, e.g. full version hosted on arXiv, HAL, or other respository/website
%%\relatedversiondetails[linktext={opt. text shown instead of the URL}, cite=DBLP:books/mk/GrayR93]{Classification (e.g. Full Version, Extended Version, Previous Version}{URL to related version} %linktext and cite are optional
%
%%\supplement{}%optional, e.g. related research data, source code, ... hosted on a repository like zenodo, figshare, GitHub, ...
%%\supplementdetails[linktext={opt. text shown instead of the URL}, cite=DBLP:books/mk/GrayR93, subcategory={Description, Subcategory}, swhid={Software Heritage Identifier}]{General Classification (e.g. Software, Dataset, Model, ...)}{URL to related version} %linktext, cite, and subcategory are optional
%
%%\funding{(Optional) general funding statement \dots}%optional, to capture a funding statement, which applies to all authors. Please enter author specific funding statements as fifth argument of the \author macro.
%
%%\acknowledgements{I want to thank \dots}%optional
%
%%\nolinenumbers %uncomment to disable line numbering
%
%%Editor-only macros:: begin (do not touch as author)%%%%%%%%%%%%%%%%%%%%%%%%%%%%%%%%%%
%\EventEditors{John Q. Open and Joan R. Access}
%\EventNoEds{2}
%\EventLongTitle{42nd Conference on Very Important Topics (CVIT 2016)}
%\EventShortTitle{CVIT 2016}
%\EventAcronym{CVIT}
%\EventYear{2016}
%\EventDate{December 24--27, 2016}
%\EventLocation{Little Whinging, United Kingdom}
%\EventLogo{}
%\SeriesVolume{42}
%\ArticleNo{23}
%%%%%%%%%%%%%%%%%%%%%%%%%%%%%%%%%%%%%%%%%%%%%%

\newcommand{\calP}{\mathcal{P}}
\newcommand{\calC}{\mathcal{C}}
\newcommand{\calV}{\mathcal{V}}
\newcommand{\calI}{\mathcal{I}}  
\newcommand{\calO}{\mathcal{O}}  
\newcommand{\calH}{\mathcal{H}}
\newcommand{\calE}{\mathcal{R}}
\newcommand{\calR}{\mathcal{R}}
\newcommand{\calD}{\mathcal{P}}
\newcommand{\cald}{\mathcal{C}}
\newcommand{\calS}{\mathcal{S}}
\newcommand{\calB}{\mathcal{B}}
\newcommand{\calT}{\mathcal{T}}
\newcommand{\calG}{\mathcal{G}}
\newcommand{\ICS}{\mathit{ICS}}
\newcommand{\cnt}{\mathit{cnt}}
\newcommand{\stake}{\mathit{power}}
\newcommand{\QC}{\mathit{QC}}
\newcommand{\self}{\mathit{self}}
\newcommand{\success}{\mathit{success}}
\newcommand{\failure}{\mathit{failure}}
\newcommand{\seed}{\mathit{seed}}
\newcommand{\rand}{\mathit{rand}}
\newcommand{\hash}{\mathit{hash}}
%\newcommand{\proof}{\mathit{proof}}
\newcommand{\VRF}{\mathit{VRF}}
\newcommand{\totalStake}{\mathit{totalPower}}
\newcommand{\power}{\mathit{power}}
\newcommand{\totalPower}{\mathit{totalPower}}
\newcommand{\economicValue}{\mathit{economicValue}}
\newcommand{\socialValue}{\mathit{socialValue}}
\newcommand{\work}{\mathit{work}}
\newcommand{\workDelta}{\mathit{workDelta}}
\newcommand{\socialDelta}{\mathit{socialValueDelta}}
\newcommand{\sk}{\mathit{sk}}
\newcommand{\pk}{\mathit{pk}}
\newcommand{\ang}[1]{$\langle \mathrm{#1} \rangle$}
\newcommand{\msg}[2][]{$\langle #2 \rangle_{#1}$}
\newcommand{\txt}[1]{\texttt{#1}}
\newcommand{\m}[1]{\mathit{#1}}
\newcommand{\msc}[1]{\mathrm{\textsc{#1}}}
\newcommand{\tsc}[1]{\textsc{#1}}
\newcommand{\msgt}[2][]{$\langle \texttt{#2} \rangle_{#1}$}
\newcommand{\angt}[1]{$\langle \mathrm{\texttt{#1}} \rangle$}
\newcommand{\twothird}{\frac{2}{3}}
%\newcommand{\Pr}{\mathit{Pr}}
\newcommand{\GST}{\mathit{GST}}
\DeclareMathOperator*{\argmax}{arg\,max}


\newcommand{\Upon}{\textbf{upon }}
\newcommand{\From}{\textbf{from }}
\newcommand{\while}{\textbf{while }}
\newcommand{\Do}{\textbf{do }}
\newcommand{\Respond}{\textbf{respond }}
\newcommand{\Start}{\textbf{start }}
\newcommand{\Broadcast}{\textbf{broadcast }}
\newcommand{\Send}{\textbf{send }}
\newcommand{\Schedule}{\textbf{schedule }}
\newcommand{\true}{\textbf{true }}
\newcommand{\false}{\textbf{false }}
\newcommand{\IndState}[1][1]{\State\hspace{#1\algorithmicindent}}
\algrenewcommand\algorithmicfor{\textbf{upon}}
\algrenewcommand\algorithmicwhile{\textbf{until}}
%\algrenewcommand\algorithmicforall{\textbf{for}}
\mathchardef\mhyphen="2D
\newcommand{\bbV}{\mathbb{V}}
\newcommand{\bbP}{\mathbb{P}}
\newcommand{\bbS}{\mathbb{S}}
\newcommand{\bbT}{\mathbb{T}}
\newcommand{\frakP}{\frak{P}}

\newtheorem{theorem}{Theorem}
\newtheorem{proposition}{Proposition}
\newtheorem{lemma}{Lemma}
\newtheorem{definition}{Definition}
\newtheorem{assumption}{Assumption}

\makeatletter
\newenvironment{breakablealgorithm}
  {% \begin{breakablealgorithm}
   \begin{center}
     \refstepcounter{algorithm}% New algorithm 
     \hrule height.8pt depth0pt \kern2pt% \@fs@pre for \@fs@ruled
     \renewcommand{\caption}[2][\relax]{% Make a new \caption
       {\raggedright\textbf{\fname@algorithm~\thealgorithm} ##2\par}%
       \ifx\relax##1\relax % #1 is \relax
         \addcontentsline{loa}{algorithm}{\protect\numberline{\thealgorithm}##2}%
       \else % #1 is not \relax
         \addcontentsline{loa}{algorithm}{\protect\numberline{\thealgorithm}##1}%
       \fi
       \kern2pt\hrule\kern2pt
     }
  }{% \end{breakablealgorithm}
     \kern2pt\hrule\relax% \@fs@post for \@fs@ruled
   \end{center}
  }
\makeatother
\setlength{\parskip}{1ex plus 0.5ex minus 0.2ex}
\setlength{\parindent}{0pt}

\newcommand{\srinidhi}[1]{\textcolor{blue}{Srinidhi: #1}}
\newcommand{\rishal}[1]{\textcolor{teal}{Rishal: #1}}

\newcommand{\codename}{\textsc{IFP}}
%%%%%%END OF PREAMBLE%%%%%%%%%%
\begin{document}
    \begin{abstract}
    \end{abstract}
    \maketitle
    %\section{Introduction}

%\subsection{Organisation}
%In \Cref{sec:contribution} we introduce our main contributions. 
%The setting of the protocol and preliminaries are described in \Cref{sec:model}.
%We provide a complete description of our protocol in \Cref{sec:protocol} followed by an analysis in \Cref{sec:analysis}. 
%It is followed by an analysis of our protocol in \Cref{sec:analysis}.
%Finally, we conclude in \Cref{sec:conclusion}.
%\newpage
%\section{Overview}
\section{Introduction}
\label{sec:contribution}

Blockchain systems starting with Bitcoin~\cite{Nakamoto09}, democratized the ability to maintain a global and shared distributed ledger of transactions without the need for a trusted third party.
Ethereum~\cite{But15} further built on this vision, building an open-access general purpose decentralised application platform - a functionality which most modern blockchains provide. 
In the modern world, human interactions are becoming increasingly digital. 
Such interactions are highly contextual - for instance, an interaction where a person pays for a coffee is different from the person buying real estate although both involve a money transfer. 
Therefore, a network for such digital human interactions must integrate their contextual nature. 
However, in today’s general purpose smart contract platforms, the state intelligence of participants for each application is siloed with the contracts themselves.
%This results in inflexibility as the largely context sensitive real world interactions cannot be seamlessly catered to based on their requirements as the intelligence is not present to do so. 
This results in inflexibility as the intelligence to seamlessly cater to the largely context-sensitive real world interactions is not available.
For instance, in the example interactions given earlier, the contextual difference between the two cannot be leveraged in traditional blockchains to incorporate more trust in the latter interaction due to the siloed nature of intelligence.
If the contextual intelligence of a participant is leveraged across all applications, a participant-centric network is created whose functionality powers the digitally interacting future for humans.
This requires a fundamental infrastructure shift, moving the intelligence from applications hosted on smart contracts to the network itself. 
This intelligence in the network is anchored to the participants, making participant preferences at the centre of applications built on the network. 
This is captured by the \emph{context object} (analogous to state) of participants, having four dimensions of Value, Logic, Files, and Trust.  

\begin{figure}[h!]
	\tikzfig{intro}
	\caption{Interactions between Alice, Bob, Charlie, and David}
	\label{fig:intro}
\end{figure}

Consider 4 distinct participants, Alice, Bob, Charlie, and David, who interact with each other in a blockchain system. 
Together, they submit two distinct transactions as given in Figure 1 - Alice sends 10 coins to Bob and Charlie sends 5 coins to David. 
In a traditional blockchain, validators that receive both these transactions ignore the context of the interacting participants and communicate with all other validators to decide on the order of the transactions. 
Here, the consensus mechanism is participant agnostic and only concerns the global state. 

In this work, we present the Interaction Finality Protocol (IFP) - a consensus protocol powering participant-centric blockchains which fundamentally integrates the participant into the consensus mechanism. 
IFP achieves consensus on updated context objects of participants involved in interactions. 
It powers a blockchain where each participant has its respective separate chains that intersect only when an \emph{interaction} (analogous to transaction) occurs between the participant. 
Here, the current context object of a participant can be found in the head of its respective chain. 
This decouples consensus, providing the ability to integrate participants’ trust preferences into the consensus mechanism. 
However, when multiple concurrently occurring interactions involve a common participant, it is crucial to ensure that they are ordered globally among all other concurrent interactions. 

To achieve this, IFP’s consensus mechanism associates a committee of validators with each participant.
Every interaction forms a unique Interaction Consensus Set (ICS) consisting of the committees of the interacting participants, along with a randomly chosen set of validators from the network that order the interaction together. 
The ICS ensures that the validators maintaining the participant’s chains order the interaction globally among all other concurrent interactions. 
The randomly chosen set of validators provides the power of the entire network, having to corrupt a significant portion of the entire network to affect IFP’s safety. 
Further, this ability to concurrently process interactions can lead to significant performance improvements in terms of throughput.


%Intro written by Srinidhi:%%%%%%
%
%Blockchain systems, starting with Bitcoin~\cite{Nakamoto09} and Ethereum~\cite{But15}, democratized the ability to maintain a global and shared distributed ledger of transactions without the need for a trusted third party. These systems rely on the ability of participants to demonstrate sufficient proof of work (computational effort) to ensure resilience against malicious actors. However, the requirements imposed by the proof-of-work mechanism constrains the scalability of these systems, making it difficult to handle a large number of transactions efficiently. 
%
%Driven by the scalability challenge, alternative consensus mechanisms that relied on traditional Byzantine fault tolerance (BFT) algorithms were developed. Systems such as Tendermint~\cite{BKM18}, Algorand~\cite{GHMVZ17}, and Hotstuff~\cite{YMRGA19} introduced proof-of-stake based consensus mechanisms motivated by traditional BFT algorithms such as PBFT~\cite{CL99} and its variants. 
%
%Both class of systems involved two distinct actors interacting with the blockchain - a participant who submits transactions and a validator who orders these transactions globally among all participants. However, the validator's core ordering mechanism (proof-of-work or proof-of-stake) is agnostic to the participant and the context of the particular transactions. 
%
%Consider 4 distinct participants, Alice, Bob, Charlie, and David, who interact with each other in a blockchain system. Together, they submit two distinct transactions as given in \Cref{fig:intro} - Alice sends 10 coins to Bob and Charlie sends 5 coins to David. In a traditional proof-of-stake blockchain, validators that receive both these transactions ignore the context of the interacting participants and communicate with all other validators to decide on the order of the transactions. However, since the transactions involves a disjoint set of participants, any ordering of the transaction will result in the same global state, an optimization that is not possible to exploit in current blockchain systems. 
%
%\begin{figure}[h!]
%\tikzfig{intro}
%\caption{Interactions between Alice, Bob, Charlie, and David}
%\label{fig:intro}
%\end{figure}
%
%In general, incorporating the context of the participants under an interaction-centric consensus mechanism can lead to significant performance improvements in terms of throughput and latency.
%
%In this work, we present \codename{}, a novel blockchain powered by a new consensus algorithm that fundamentally integrates the different actors into the decision mechanism. \codename{} is a proof-of-stake based blockchain where each participant has its respective separate chains that intersect only when an interaction occurs between the participants. This decouples consensus between different participants, allowing for multiple interactions among different participants to be agreed upon simultaneously.
%
%However, when multiple interactions, occurring concurrently, involve a common participant, ensuring that the interactions are ordered globally among all other concurrent interactions is crucial. To achieve this, \codename{}'s consensus mechanism associates with each participant a context which is a set of validators that are associated with the participant. 
%% When a participant interacts with other participants, their respective contexts ensure that the interaction is ordered globally among all other concurrent interactions.
%Every interaction forms a unique \emph{Interaction Consensus Set} (\emph{ICS}) consisting of the contexts of the interacting participants, along with a randomly chosen set of validators from the network that order the interaction together. The ICS ensures that the validators maintaining the participant's chains order the interaction globally among all other concurrent interactions.
%
%Having a public and small context for each participant makes \codename{} more prone to bribery and collusion attacks. Validators that are part of a participant's context will order transactions preferentially in favour of the participant for a fee. A key feature of \codename{} that mitigates bribery and collusion attacks is its ability to update the context of a participant dynamically. Specifically, the context of each participant is stored in a special chain where transactions on this participant updates the context.
%%%%%%%%%%%%%%%%%%%

%\srinidhi{I'm not happy with this introduction. It lacks a concrete contribution and is very handwavy. We say we can gain something but don't back it up. We say here is a cool new algorithm but don't say what we gain from it. I'll take another pass at this later if we can position it better. Even writing it feels very shallow.}

% This is a draft story to motivate and position this work. Will fill it with concrete text later
% \todo[inline]{
	% - Traditional blockchain systems evolved from (expert only) running your own node to mine currency to a general consumer driven service driving financial inclusion
	% - To bridge the gap, two distinct roles emerged in modern blockchain systems - a participant and a validator
	% - A participant engages with the blockchain to submit transactions that exchange currency and other assets (in the form of smart contracts)
	% - A validator/node accepts transactions and orders them globally among all participants
	
	% - In order to ensure total ordering of transaction, validators typically relied on a proof of work scheme. Validators play a computation game and the node that wins the game decides the order of the next set of transactions (blocks).
	% - However this does not scale.
	% - Alternatively, proof of stake schemes involve a selected validator proposing a set of transactions and use a standard off the shelf consensus algorithm to ensure all other validators order transactions as proposed.
	% - However, the core consensus mechanism is agnostic to the different actors - participant and validator - that engage with the blockchain
	
	% - In this work, we present \codename{}, a novel blockchain powered by a new consensus algorithm that fundamentally integrates the different actors into the decision mechanism
	% - \codename{} is a proof of stake based blockchain where different participants their respective separate chains that only interact with each other if the transaction involves them.
	% - <Example diagram and an accompanying text>
	
	% - \srinidhi{What do we concretely gain? Can we give guarantees about that and later evaluate? - Performance in terms of throughput of the transactions? If not, we can just present the algorithm but we should be saying something about asymptotic gains then. Essentially the question is - What is our contribution with the new algorithm? Other than its unique and different.}
	
	% - A short explanation of normal mode working and motivating the different terms
	% - <Here we should introduce and intuitively explain the DAG and Interaction-Centric Consensus>
	% }


% To ensure total ordering of interactions, Bitcoin relied on a proof of work scheme. 
% Here, validators solve computationally difficult puzzles and validators that solve this puzzle decide the next set of interactions (called a block) to be committed in the total order. 
% However, this approach has significant computational requirements due to the inherent computational demand for solving the puzzles and moreover does not scale for high throughput of interactions. 
% An alternative mechanism involves a set of validators being selected and a BFT consensus algorithm being run within this committee that decides the order of blocks. 
% Here, the core consensus mechanism is agnostic to the participants that engage with the blockchain. 
% This provides better throughput but does not scale with increase in number of validators in the network. 
% In these traditional approaches, interactions performed between participants are batched into blocks and agreed upon, forming a chain of blocks, each of whom represent the network's global state. 
% In this work, we present \codename{}, a novel blockchain powered by a new consensus algorithm that fundamentally integrates participants into the decision mechanism. 
% \codename{} is a blockchain where each participant has its respective separate chains that intersect only when an interaction occurs between the participants. 
% This decouples consensus between different participants, allowing for multiple interactions among different participants to be agreed upon simultaneously. 
% These interactions between participants are committed only to the chains of the participants involved. 
% The execution of an interaction updates the state of the interacting participants and this is represented by a \emph{tesseract}. 
% Thus, the latest state of a participant is found in the latest committed tesseract with the participant. 
% The global state of the network is obtained by simply taking the latest states of all the participants. 
% The chains of each participant intersect at tesseracts involving them to form a log of the global state called the \emph{Multi-Linked DAG} or simply \emph{mDAG}.
% This consensus procedure involves validators associated with participants called its \emph{context}. 
% The \emph{context} of a validator captures various participant attributes in the Krama blockchain that we do not describe here in addition to a set of associated validators.
% When performing interactions, the sender of the interaction can include a set of validators called the \emph{preferential context} that are to be included in achieving consensus on the interaction. 
% This preferential set allows participants to include its trusted validators in the interaction's consensus providing flexibility to participants by making each consensus procedure tailored for the interaction.
% This flexibility is a direct consequence of the mDAG's design, having chains for each participant that intersect only when interactions are made. 
% For each interaction, an \emph{Interaction Consensus Set} (\emph{ICS}) is formed, consisting of the contexts of the interacting participants, the preferential set if mentioned, and a randomly sampled set of stochastic validators.
% This stochastic set leverages the power of the network's size, preventing collusion among the relatively small and static participant contexts. 
% Validators in an ICS for an interaction can simultaneously participate in consensus for other interactions with disjoint participant sets.
% A leader based consensus algorithm with three round-trips of communication driven by reception of quorums of messages is run by validators in the ICS. 
% The duration of time for this consensus procedure is predetermined, at the end of which validators terminate, acting as a timeout-driven view-change mechanism. 
% This ensures that validators are available to enter consensus instances for interactions in-sync. 

% The flexibility obtained by having a separate ICS for each interaction allows for a participant-driven network where participants' preferences can be incorporated into the consensus mechanism.


\subsection{Organisation}

In what follows, we first present the preliminary definitions in~\cref{sec:preliminaries} and the system model in~\cref{sec:model}.
We then present the core \codename{} protocol in \Cref{sec:protocol} followed by its analysis in \Cref{sec:analysis}. Finally, we discuss related work in \Cref{sec:relatedwork} and conclude in \Cref{sec:conclusion}.

    \section{Preliminaries}\label{sec:preliminaries}

Two actors drive the algorithm in \codename{}. First, validators (denoted by $\bbV$), that run the protocol to reach consensus. Second, Participants ($\bbP$) that submit \textit{interactions} between each other to the validators. Each participant $p \in \bbP$ defines a fixed-size set of validators $C_p \subseteq \bbV$ ($|C_p| = n_c$) that mandatorily participate when achieving consensus for interactions involving the participant $p$. We will refer to this set of validators as its \textit{context}. The context of each participant is public and changes periodically.

% In this section, we establish some definitions that are used in the rest of the paper.

% %\subsection{Preliminaries}

% %\subsubsection{Participant}\label{subsubsec:participant}
% We have two different types of actors: \emph{validators} and \emph{participants}. 
% Validators run the Krama protocol and reach consensus on interactions made between participants.
% Let the sets of validators and participants in the network be denoted as $\bbV$ and $\bbP$, respectively. 
% Each participant $p\in \bbP$ is assigned a fixed-size set of validators $\calC_p\subseteq \bbV$ called its \emph{context} that participate in achieving consensus on interactions involving the participant. 
% The size of the context $|\calC_p|$ is denoted as $n_\calC$. 
% All validators are aware of the context of each participant. 
% Additionally, the protocol ensures that the context is updated periodically. \todo{'periodically' instead of 'epoch' because epoch not yet def.} 

The state of each participant (capturing the assets of the participant) is stored by each validator. 
Formally, we define the state of each participant as follows. 

\paragraph{\textbf{State}}
%Each participant $p\in\bbP$ has a \emph{state} $S_p$ updated through interactions.
The \emph{state} of a participant $p$ is defined as $S_p = (l_p, \calC_p)$ where $l_p$ is the valuation of its local variables and $\calC_p$ its context as defined earlier. 
Subsequently, $\bbS_p$ denotes the set of possible states for a participant $p$.
We define the \emph{global state} of the network as a union of the state of all participants, $S = \bigcup_{p\in\bbP} S_p$. and $\bbS$ is the set of all possible global states.


%At the start of the protocol, the participant states are all initialised as the genesis state.%, con with the global state containing information of the 


%Each validator $v\in\bbV$ maintains a copy of the global state $S^v$, called its \emph{history}. 
%Interactions performed by a participant update its participant state, represented by a tesseract, resulting in a chain of tesseracts called its \emph{participant chain}.
%The number of tesseracts present in a participant $p$'s chain is called its \emph{height} and is denoted by $h_p$.

%Thus, interactions performed by a participant results in a chain of states called its \emph{participant chain}. 
%Each participant's chain contains a public seed in the latest tesseract that is updated when a new tesseract is committed.

\paragraph{\textbf{Interaction}}
The state of each participant (and consequently the global state) is updated through \emph{interactions} (analogous to transactions).
Participants send interactions to the network which are processed by the validators.
%Old - 26 May 2025
% An interaction $I$ between a set of participants $P_I$ consists of 
% \begin{enumerate}
%     \item The \emph{action} $A_I$, \srinidhi{unclear what you mean here} transitioning the states of the interacting participants. It is given as a set of transition functions $A_I^p:\bbS_p\to\bbS_p$ for each $p\in P_I$. 
%     \item Optionally, a set of validators $\calP_I$ called \emph{preferential context} \srinidhi{You need to motivate this earlier. Maybe it would make sense to restructure such that the protocol comes first.} \rishal{Yes. But might be difficult to go over the protocol without these definitions in place. Can instead move this to the interaction consensus set defn and say that Additionally, interactions include the preferential context given by the participants...} that are to be included in the consensus process for the interaction.
% \end{enumerate}
% \srinidhi{Validators collectively update the global state $S$ by executing an interaction as follows, ...}
% The execution of an interaction $I$ with current global state $S$ is given by the global state transition function $\delta:\bbS\times\mathbb{I}\to \bbS$ and results in an updated global state $S'=\delta(S, I)$. 
% Here, $\mathbb{I}$ is the set of possible interactions. \srinidhi{Move this to the definition of interactions}
% This state transition $\delta$ is defined as follows: 
% \begin{enumerate}
%     \item For each $p\in P_I$, $S^{\prime}_p=A_I^p(S_p)$ is the participant state obtained on applying the action $A_I$ on $p$, and  
%     \item For each $q\notin P_I$, $S'_q=S_q$. 
% \end{enumerate}

%New - 26 May 2025
Formally, we define an interaction $I$ between a set of participants $P_I$ as a set of transition functions $\delta_I^p:\bbS_p\to\bbS_p$ for each $p\in P_I$. An interaction abstracts the semantics of the underlying transitions between the states of the participants involved. Naturally, each validator updates the state of all participants $p \in P_I$ by applying the corresponding transition function $S'_p = \delta_I^p(S_p)$.

% For the purpose of simplicity, we assume that an interaction consists of only two participants, i.e., $|P_I|=2$. However, the abstract definition of an interaction allows the protocol to be extended to any number of participants.

% The set of participants in an interaction $I$ is denoted as $P_I$.
% Validators update the participant states $S_p$ for each interacting participant $p\in P_I$ by executing the interaction as follows. 
% For a participant $p\in P_I$, the execution of the interaction results in the updated state $S'_p=\delta_I^p(S_p)$ obtained on application of the corresponding transition function present in the interaction. 
% Thus, the execution of an interaction results in a set of updated participant states for the interacting participants. 
% This is represented by a \emph{tesseract} and is described in detail below. 

\paragraph{\textbf{Tesseract}}
From a given global state $S$ and an interaction $I$, a \textit{Tesseract} $T$ represents the set of updated states for each participant involved in the interaction. Formally, $T=\{\delta_I^p(S_p)\mid p\in P_I\}$ and $\bbT$ is the set of all possible Tesseracts. Validators internally log the Tesseract for each interaction and the log serves as snapshots of the global state. The log is stored using a data structure called Multi-Linked DAG.

% \srinidhi{It wouldn't hurt to use known terminology just to help the reader understand. By this point, we've already seen many new terms and the reader is left wondering what maps to what.}
% A \emph{tesseract} represents the updated state of the interacting participants after executing an interaction, serving as a snapshot of the participant states. 
% The execution of an interaction $I$ with global state $S$ results in the tesseract $T=\{\delta_I^p(S_p)\mid p\in P_I\}$ consisting of the updated participant states .
% The tesseracts resulting from execution of interactions by validators result in a growing log of tesseracts called the \emph{multi-linked DAG}.

\begin{figure}
\tikzfig{mdag}
\caption{The Multi-Linked DAG. The edge colours represent participants.}
\label{fig:mdag}
\end{figure}

\paragraph{\textbf{Multi-Linked DAG}}
% All validators store the \emph{multi-linked directed acyclic graph} (\emph{mDAG}) locally and is referred to as their \emph{history}. 
%The execution of interactions result in the creation of tesseracts which are organised into a growing log of tesseracts in the network. 
%This is organised as a \emph{multi-linked directed acyclic graph} (\emph{mDAG}).
%Each validator stores this mDAG locally and is referred to as its \emph{history}. 
% Tesseracts representing these states are linked to tesseracts with previous states for each interacting participant through directed edges to the previous state labeled by the participant. 
% This graph made with nodes as tesseracts and directed edges labelled by participants forms a DAG and is referred to as the \emph{multi-linked composite DAG} or simply \emph{mDAG}. 
Multi-Linked DAG (mDAG) of a validator $v$ is a labelled directed graph $\calG_v=(\mathbb{T}, E, \ell)$. The nodes of the graph $\bbT$ are tesseracts and the edges $E \subseteq \bbT \times \bbT$ are labeled by participants (given by $\ell: E \to \bbP$). A Multi-Linked DAG is given in \Cref{fig:mdag}.

% with the set of labels $\bbP$ given by $\calG^v=(\mathbb{T}, E, \ell)$ where $\mathbb{T}$ is the set of tesseracts in the history that act as vertices, $E\subseteq \mathbb{T}\times\mathbb{T}$ is the set of directed edges given as a tuple of vertices, and $\ell:E\to\bbP$ is the edge labeling function that associates each edge with a participant. 

\paragraph{\textbf{Tesseract Relationships}}
The Tesseracts form intuitive relationships driven by the edges of the mDAG. We say $T$ is a \textit{$p$-parent} of $T'$ if there exists an edge $e=(T', T) \in E$ with the label $l(e)=p$. Conversely, $T'$ is the \textit{$p$-child} of $T$, and $T\preceq_p T'$ is the transitive closure of the \textit{parent} relationship ($T$ is a \textit{$p$-ancestor} of $T'$).

Intuitively, the parent-child relationship defines a partial order on the Tesseracts. The order is crucial to break ties between two competing Tesseracts which are updating the states of a common participant $p$. Formally, we say $T$ conflicts with $T'$ for a given participant $p \in P_T \cap P_{T'}$ if neither is a \textit{$p$-ancestor} to the other.



% We introduce some notions of relationships between tesseracts in the mDAG that are used in the rest of the paper. 
% Let $T$ and $T'$ be tesseracts present in the history of a validator.
% The tesseract $T$ is said to be a \emph{$p$-parent} of $T'$ if there exists an edge $e=(T', T)\in E$ where $\ell(e)=p$ in the mDAG. Here, $T'$ is also called a $p$-child of $T$.
% The tesseract $T$ is said to be a \emph{$p$-ancestor} to $T'$, denoted by $T\preceq_p T'$, if there exists a path labeled by $p$ from $T'$ to $T$ in the mDAG.
% The tesseracts $T$ and $T'$ are said to be \emph{conflicting} if there exists $p\in P_{T}\cap P_{T'}$\todo{notation unused} such that neither is a $p$-ancestor to the other. 

%\srinidhi{Motivate why we need to define agree, how can it possibly diverge?}
A tesseract $T$ corresponding to interaction $I$ is said to \emph{extend} a validator $v$'s state $S$ if $T=\{\delta_I^p(S_p)\mid p\in P_I\}$.
A tesseract is \emph{committed} to the history of a validator by linking it with directed edges labeled by the participant to the tesseracts with the participant states of each of the interacting participants. 
More precisely, a tesseract $T$ is committed to the history $\calG_v=(\bbT, E, \ell)$ of validator $v$, resulting in an updated history $\calG'_v=(\bbT', E', \ell')$ where %(1) $\bbT'=\bbT\cup \{T\}$, (2) $E' = E\cup \{(T, T_p)\mid p\in P_T\}$, and (3) $\ell'=\ell\cup\{((T,T_p), p)\mid p\in P_T\}$ 
\begin{enumerate}
    \item  $\bbT'=\bbT\cup \{T\}$, 
    \item  $E' = E\cup \{(T, T_p)\mid p\in P_T\}$, and
    \item $\ell'=\ell\cup\{((T,T_p), p)\mid p\in P_T\}$.
\end{enumerate}
Here, for each $p\in P_T$, $T_p$ is the tesseract with the participant state $S_p$ in $\calG_v$. 
Thus, for each interacting participant $p$, the tesseracts with the earlier participant states are the $p$-parents of the committed tesseract. 


% \paragraph{\textbf{State Retrieval}}
% The latest state of a participant $p$ is found in the Tesseract $T_p$ that has no $p$-children. As mentioned earlier, the global state $S$ of the system is a composition of the latest Tesseracts $T_p$ for all participants $p$.

% The current state of a participant is found in its latest state transition's tesseract.
% Similarly, the current global state is the composition of all such participant states. 
% More precisely, for a participant $p$, its current state $S_p$ in the history $\calG=(\bbT, E,\ell)$ is found in the tesseract $T_p\in\bbT$ having a $p$-parent and no $p$-children. 
% Similarly, the current global state $S$ in the history is found in the set of such tesseracts $\{T_p\mid p\in \bbP\}$.


\paragraph{\textbf{Height}}
For a given Tesseract $T$, there exists many paths from it to the genesis Tesseract $G$ in the global state. However, only one unique path for a given participant $p$ exists i.e. the path with edges labelled $p$. Therefore, the height of a Tessetact is defined as a function of the participant. Formally, $T.\m{height}(p)$ is the length of the path following the edges labeled by $p$ from the Tesseract $T$ to the genesis Tesseract $G$.

% In a validator's history, a tesseract can have more than one path from it to the genesis tesseract. 
% However, it has a unique path labeled by a given interacting participant.
%Thus, the notion of height of a tesseract is restricted to that of its 
%In Krama, the notion of the height of a tesseract is restricted to that of its interacting participants; there is no notion of a tesseract's height as it can have more than one path to the genesis tesseract. 
% The \emph{height} of a tesseract for an interacting participant is the length of the path labeled by the participant from the tesseract to genesis. 
% More precisely, the height of tesseract $T$ for an interacting participant $p$ is given by the function $T.\m{height}:P_T\to\mathbb{N}$, where $T.\m{height}(p)$ for $p\in P_T$ is the height of tesseract $T$ for participant $p$.
% At genesis, the height of all participants is zero, that is, $G.\m{height}(p)=0$ for all $p\in \bbP$ where $G$ is the genesis tesseract. 
% The height of a tesseract $T$ for a participant $p$ is the height of its $p$-parent incremented by one, that is, $T.\m{height}(p)=T'.\m{height}(p)+1$ where $T'$ is the $p$-parent of $T$.
%When a tesseract is committed, the height of each interacting participant is incremented by one.



% \subsubsection{mDAG Properties}
% Since the mDAG is built by committing new tesseracts as given above, the history $\calG=(\bbT, E, \ell)$ stored by a validator has the following properties, 
% \begin{enumerate}
%     \item If an edge $e=(T_2, T_1)$ is present in $E$ with $\ell(e)=p$ then $p\in P_{T_1}$ and $p\in P_{T_2}$ and further, $T_2$ was committed after $T_1$,
%     \item If a tesseract $T$ is present in $\bbT$, then there exist edges $e_p$ for each $p\in P_T$ with labels $\ell(e_p)=p$ to the tesseract with the earlier participant state when $T$ was committed,
%     \item For each tesseract $T\in \bbT$, $T$ has in-edges and out-edges labeled by $p$ if and only if $p\in P_T$.
% \end{enumerate}


\paragraph{\textbf{Interaction Consensus Set}}
When a new interaction $I$ is proposed, a set of validators communicate to reach consensus on the interaction. We define the set of validators as \emph{Interaction Consensus Set} (ICS) for the interaction $I$ and is denoted as $\mathit{ICS}(I)$.

The ICS consists of,
\begin{enumerate}
    \item The context of each interacting participant $p\in P_I$, denoted as $\calC_p$.
    \item A \emph{stochastic set} of validators randomly sampled from the set of validators, denoted as $\calS$.
    \item The \emph{initiating} participant of the Interaction is allowed to specify a set of validators called the \emph{preferential context} denoted as $\calP_I$ that are to be included in the consensus process for the interaction.
\end{enumerate}

The context of each participant $p \in P_I$ is sufficient to ensure that the protocol extends the latest state of each participant. However, the set is fixed and predetermined, and hence allows for bribery attacks \cite{KKZ24}. Therefore, we include as a second component a variable \emph{stochastic set} of validators randomly sampled from the set of validators. This ensures that the interaction is processed by a random subset of validators, preventing collusion and bribery attacks. 


% \subsubsection{Interaction Consensus Set}
% To reach consensus on an interaction $I$, the validators from the context $\calC_p$ of each interacting participant $p\in P_I$ are required to be present. 
% Providing flexibility to the participant, the sender of an interaction $I$ can additionally provide a set of validators called the \emph{preferential context} 
% denoted as $\calP_I$ that are to be included in the consensus process for the interaction. 
% This incorporates the participant's trust preferences in the consensus mechanism.
% Now, the \emph{context} of an interaction $I$, given as $\msc{ctx}(I)=\bigcup_{p\in P_I} \calC_p \cup \calP_I$ comprises of the interacting participants' contexts as well as the preferential context of the interaction. 
% However this is vulnerable to collusion and bribery attacks as the relatively static contexts are only updated periodically.
% Thus, to reach consensus, an \emph{interaction consensus set} (or simply, \emph{ICS}) consisting of (1) the context of each interacting participant, (2) a \emph{stochastic set} of validators randomly sampled from the set of validators, and, if mentioned, (3) the preferential context present in the interaction. 
% The stochastic set is typically denoted by $\calS$ and is the same size as that of the participants' contexts ($n_\calC$). \todo{is size relevant?}
% The ICS for an interaction $I$ is thus given by $\mathit{ICS}(I) = \bigcup_{p\in P_I} \calC_p \cup \calS \cup \calP_I$.  
%In a given view, honest validators cannot engage in consensus for multiple interactions with a common participant.
%This prevents conflicting tesseracts being committed in the common participant's chain. \srinidhi{unclear what you mean}

% \begin{example}
%     Consider two interactions $\calI_{P}$ and $\calI_{P'}$ where $P = \{p_1, p_2\}$ and $P' = \{p_3, p_2\}$. 
%     An honest validator cannot engage in consensus for both $\calI_P$ and $\calI_{P'}$ in a given view.
% \end{example}

\paragraph{\textbf{Quorum}}

A \emph{quorum} is a set of validators that can reach consensus on an interaction. In the context of our protocol, a quorum is essential for ensuring that the interaction is agreed upon by a sufficient number of validators, thereby providing resilience against Byzantine faults. 

In particular, a quorum consists of supermajorities from (1) the context of each interacting participant, (2) the stochastic set, and (3) the preferential context if present where a \emph{supermajority} is a subset of $2f+1$ unique validators from the set of validators $X$ of size $n=3f+1$.


% \subsubsection{Supermajority and Quorum}

% Let $X$ be a set of validators of size $n=3f+1$.
% For simplicity, let each subset of validators $V\subseteq\bbV$ considered in the rest of the paper be of size $s$ such that $s=3b+1$ for some positive integer $b$.
% A \emph{supermajority} over the set $X$ is a subset of $2f+1$ unique validators. 
% A \emph{quorum} over an ICS consists of supermajorities from (1) the context of each interacting participant, (2)  the stochastic set and (3) the preferential context if present. 
% Similarly, a quorum over an interaction's context consists of only (1) and (3) above.
%A \emph{quorum certificate} is obtained from a quorum of votes. 
%A \emph{minority certificate} from an ICS is a collection of votes from one of the sets $X$ constituting the ICS with at least $f+1$ validators and a cumulative context power of $b+1$.
%To prevent Sybil attacks in the permissionless setting, each validator is assigned a weight. 
%The actual slightly differs from the above, with each validator having a weight to prevent Sybil attacks in the permissionless setting
%Therefore, a supermajority is instead a collection of $2f+1$ voting shares. 

\section{Model}
\label{sec:model}

We now describe the setting and the assumptions under which the protocol operates. 

\subsection{Network Model}
\subsubsection{Partial Synchrony}
We model network communication as the partially synchronous model (introduced in \cite{DLS88}) with a known bound $\Delta$ on message delay and an unknown instant called Global Stabilisation Time (GST), such that every message sent by an honest validator at time $t$ to another honest validator arrives within time $\max (t, \m{GST})+\Delta$. 
 
% \subsubsection{Quasi-Permissioned}
% A validator is represented by its digital key. In Krama, time is divided into epochs of fixed length (say 24 hours).
% Our protocol works in the \emph{quasi-permissioned} setting, i.e., an individual can create a new digital key pair and the validator with the digital key pair can join the protocol at the start of the next epoch subject to a successful identity verification linking the validator to the individual. 
% Note that a single individual can have multiple validators participating in the protocol. 

% \subsubsection{Quasi-Permissionless}

% A validator is represented by its digital key. 
% In Krama, time is divided into fixed-length \emph{epochs}. % of duration, say, 24 hours.

% There is a set of validators running the protocol which is updated every epoch. 
% At any given moment, individuals can create a new digital key pair and the validator with the digital key pair can join the protocol at the start of the next epoch. 
% This is captured by the \emph{quasi-permissionless} setting given in \cite{LR24}.

% \subsubsection{Trusted Setup}
% We assume the presence of an initial participant called \emph{sarga} (Sanskrit for `genesis') with a randomly sampled context set. % generated seed present in its initial tesseract. 
% This participant's context is used for initialisation of newly joined participants. 
%\todo[inline]{How is the seed randomly generated? - Algorand does distributed key generation.\\ How is initial context selected?}

\subsubsection{Randomness}
We assume the presence of a \emph{distributed randomness beacon} that periodically provides random outputs to all validators. 

\subsection{Cryptographic Primitives}

\subsubsection{PKI and Digital Signatures}
% Each validator has a public and secret key pair.
% Krama uses BLS signatures introduced in \cite{BLS01} to sign messages. 
% We assume that every message in the protocol is signed by the sender and that digital signatures are ideal. 
% A message $m$ sent by a validator $v$ is written as \msg[v]{m}. 
% We do not mention the sender when it is clear from the context and instead simply write \msg{m}.
%%% Less Verbose: %%%
Each validator has a public-secret key pair. 
\codename{} uses BLS signatures \cite{BLS01} with ideal security assumptions. 
All protocol messages are signed by their sender. 
We denote a message $m$ from validator $v$ as \msg[v]{m}, or simply \msg{m} when the sender is clear from context.

 \subsubsection{Verifiable Random Functions}
% Verifiable random functions (VRFs), introduced in \cite{MVR99}, take a publicly available random seed and a validator's private key as input and generate a pseudo-random number and a corresponding proof as the output. 
% Any validator can verify that the pseudo-random number generated by a validator corresponds to the random seed using this proof and the public key of the validator that generated the random output.
% The VRF provides two crucial properties: (1) Pseudo-randomness---the generated number is indistinguishable from a uniformly random string by an efficient adversary when the secret key is unknown, and (2) Verifiability---with the proof given by the VRF, anyone can verify that the pseudo-random number was generated by the validator with the public key for the given seed. 
% The VRF output is deterministic for a given seed and secret key. 
% Hence, Byzantine validators cannot bias the randomness.
%%% Succinct: %%%
Verifiable random functions (VRFs) \cite{MVR99} take a public random seed and validator's secret key as input, generating a pseudo-random number and corresponding proof. 
Any validator can verify that the pseudo-random number corresponds to the random seed using the proof and the generator's public key.
VRFs provide two key properties: (1) Pseudo-randomness---the output is indistinguishable from a uniformly random string to efficient adversaries without the secret key, and (2) Verifiability---anyone can verify that the pseudo-random number was generated by the validator holding the public key for the given seed using the proof given by the VRF.
The VRF output is deterministic for a given seed and secret key.
Hence, Byzantine validators cannot bias the randomness. 

% \todo[inline]{ given (1) the publicly available seed is random and cannot be controlled by the adversary and (2) the validator cannot generate new key pairs to provide him an advantage.\\Is condition (2) necessary? I read someplace that despite this VRFs should not give him an advantage.}


\subsection{Adversary}
% A validator is \emph{honest} if it performs tasks in accordance with the protocol and can send and receive messages according to the network model. 
% A \emph{Byzantine} validator, however, can deviate from the protocol arbitrarily.
% %We consider a probabilistic polynomial time \emph{adversary} that can corrupt validators, i.e., make them Byzantine. 
% We assume that the number of Byzantine validators that can be present at any point is at most one-fifth (i.e. 20\%) of the total number of validators. 
% This, along with our network and participant context size assumptions, ensures that with high probability, the participants' contexts have at most one-third Byzantine validators.
% We assume the presence of an adversary that can perfectly coordinate the Byzantine nodes in any manner, including sending inconsistent messages and remaining silent. 
% However, the adversary cannot forge digital signatures of honest validators with non-negligible probability or interfere with the guarantees of message delivery. 
%%%% Less Verbose: %%%%%
A validator is \emph{honest} if it follows the protocol and can send and receive messages as per the network model. 
A \emph{Byzantine} validator deviates arbitrarily from the protocol but cannot forge honest validators' digital signatures or violate message delivery guarantees.
We assume that every participant context and stochastic set with $n=3f+1$ validators has at most $f$ Byzantine validators. 
% \todo[inline]{
% We assume at most 20\% of validators are Byzantine at any time. 
% Combined with our network and participant context size assumptions, this ensures participants' contexts and stochastic sets contain less than one-third Byzantine validators with high probability.
% An adversary can perfectly coordinate Byzantine nodes, including sending inconsistent messages or remaining silent, but cannot forge honest validators' digital signatures or violate message delivery guarantees.
% }
    \section{The Protocol}
\label{sec:protocol}

\codename{} protocol proceeds in a sequence of \emph{epochs}, each lasting for a fixed duration. Within each epoch, validators operate in a \emph{normal mode} where they reach consensus on Tesseracts containing the latest state of interacting participants.
Between epochs, validators enter a \emph{reconfiguration mode} to update the validator list and prepare for the next epoch.

In this section, we first describe the normal operation of the protocol, detailing how validators reach consensus on Tesseracts within an epoch. Subsequently, we outline the reconfiguration mode, which occurs between epochs to update the validator list and prepare for the next epoch. At the begining of each epoch, the validators are ordered globally and synchronized to the same view.

\subsection{Normal mode}

Each epoch is divided into a sequence of fixed-length \emph{views}, where each view lasts for a duration of $\m{viewLength}$. Validators determine their current view using their local clock, referencing the time of the previous epoch transition. Validators receive interactions from participants and attempty to reach consensus on Tesseracts within each view.

Unlike other consensus protocols, \codename{} allows for multiple interactions to be executed and agreed upon within a single view, provided that the participant sets of these interactions are mutually disjoint. This means that validators can reach consensus on multiple Tesseracts in parallel, each involving different participants i.e. have multiple \emph{consensus instances} running simultaneously. Note that executing multiple consensus instances in parallel is different from batching them together as a single block, as is done in traditional blockchain protocols. In \codename{}, each interaction results in a Tesseract that is agreed upon and committed independently.

\subsubsection{Operators}

Each \textit{consensus instance}, involving an interaction $I$ is driven by a designated \emph{operator} $O_I$. \codename{} picks operators deterministically, thereby ensuring that all validators agree on the same operator. For example, one method would be to choose the operator from the context of the participants in a round-robin manner with a deterministic tie breaker among the participants.

Once an operator is chosen, each \textit{consensus instance} proceeds in a sequence of four stages: \txt{Prepare}, \txt{Propose}, \txt{Vote}, and \txt{Commit}.

\begin{algorithm}[hptb]
\caption{The Protocol: Stage 1 - Prepare}\label{alg:step1}
\begin{algorithmic}[1]         

%\LComment{\textbf{\color{black}As an operator for the view:}}
%\For{$\m{view}.\m{begin}$}
%    \If{$\m{self} = \msc{operator}(\m{view})$}
%        %\State Select $\calI_P$ from pool of interactions such that $\m{self} =\m{operator}(\calI_P, \m{view})$
%        \State $\m{I}_v \gets \m{ixn}$ where $\m{self} =\msc{operator}(P_\m{ixn}, \m{view})$
%        \State \Broadcast \msg[]{\txt{Prepare}, I_v} to $\msc{ctx}(\m{ixn}_v)$
%    \EndIf
%\EndFor
%\item[] 
\LComment{\textbf{\color{black}As an operator:}}
\For{$\m{view}.\m{begin}$}
	\If{$\m{self} = o_I$ for interaction $I$}
		\State $\m{I}_v \gets \m{I}$
		\State \Broadcast \msg[]{\txt{Prepare}, I_v} to $\msc{ctx}(\m{I}_v)$
	\EndIf
\EndFor
\item[]	

\LComment{\textbf{\color{black}As a validator in a participant $p$'s context:}}
\For{$\m{view}.\m{begin}$}
    \State $\m{prepares}^p \gets \{\}$
\EndFor
\item[]

\While{$\m{prepareTimeout}$ \textbf{upon} \msg[\m{op}]{\txt{Prepare}, I}}
    \If{$\m{op} = o_I \land p\in P_I$}
        \State $\m{prepares}^p\gets \m{prepares^p}\cup$ \{\msg[\m{op}]{\txt{Prepare}, I}\}
    \EndIf    
\EndWhile
\item[]

\For{$\m{prepareTimeout}$}
    \State \msg[\m{op}]{\txt{Prepare}, I} $= \Call{Choose}{\m{prepares}_p}$
    %\State $\m{lockedQCs}\gets \{\m{lockedQC_v^p}\mid p\in P_I\}$
    \State \Send \msg[]{\txt{Prepare}, \{\m{lockedQC_v^p}\mid p\in P_I\}} to $\m{op}$
    \ForAll{\msg[\m{op}']{\txt{Prepare}, I'} $\in \m{prepares}^p\setminus\{\Call{Choose}{\m{prepares}^p}\}$}
        \State \Send \msg[]{\txt{PrepareNil}, \m{prepares}^p} to $\m{op}'$
    \EndFor
\EndFor

% \LComment{\textbf{\color{black}As a validator in the interaction's context:}}
% \While{$\m{prepareTimeout}$}
%     \ForAll{$p\in P_\m{self}$}
%         \State $\m{prepares}_p \gets \{\}$
%     \EndFor
%     \For{\msg[\m{op}]{\txt{Prepare}, I}}
%         \If{$\m{self} \in \msc{ctx}(I) \land \m{op} = \msc{operator}(P_I, \m{view})$}
%             \ForAll{$p\in P_I\cap P_\m{self}$}
%                 \State $\m{prepares}_p\gets \m{prepares_p}\cup$ \{\msg[\m{op}]{\txt{Prepare}, I}\}
%             \EndFor
%         \EndIf    
%     \EndFor
% \EndWhile
% \item[]

% \For{$\m{prepareTimeout}$}
%     \State $\m{chosenPrepares}\gets\{\}$
%     \ForAll{$p\in P_\m{self}$}
%         \State $\m{prepare}\gets \Call{Choose}{\m{prepares}_p}$
%         \State $\m{chosenPrepares}\gets \m{chosenPrepares}\cup \{\m{prepare}\}$
%     \EndFor
%     \ForAll{\msg[\m{op}]{\txt{Prepare}, I} $\in \m{chosenPrepares}$}    
%         \State $\m{lockedQCs}\gets \{\m{lockedQC_v^p}\mid p\in P_I\cap P_\m{self}\}$
%         \State \Send \msg[]{\txt{Prepare}, \m{lockedQCs}} to $\m{op}$
%     \EndFor
%     \ForAll{\msg[\m{op}]{\txt{Prepare}, I} $\in \bigcup\limits_{p\in P_\m{self}} \m{prepares_p}\setminus \m{chosenPrepares}$}
%         \State $\m{preparesReceived} \gets \{\m{prepares}_p\mid p\in P_I\cap P_\m{self}\}$
%         \State \Send \msg[]{\txt{PrepareNil}, \m{preparesReceived}} to $\m{op}$
%     \EndFor
% \EndFor

%\item[]
% \LComment{\textbf{\color{black}Procedures:}}
% \Procedure{Choose}{$\m{prepares}$}
%     \State $\m{operators}\gets\{\}$
%     \ForAll{\msg[\m{op}]{\txt{Prepare}, I}$\in \m{prepares}$}
%         \State 
%     \EndFor
%     \Return $\m{prepare}$
% \EndProcedure
\end{algorithmic}
\end{algorithm}

\subsubsection{Prepare}

As defined in~\cref{alg:step1}, the operator initiates the \txt{Prepare} phase by sending a \txt{Prepare} message to the \emph{interaction's context}. This message contains the interaction and serves as a discovery mechanism for validators in the context to respond with their latest state information. We refer to the participant contexts of an interaction $I$ together as the \emph{interaction's context} denoted $\msc{ctx}(I)$. In the \txt{Prepare} phase, each validator (in the context of a participant $p$) collects \txt{Prepare} messages from operators and until a timeout is reached, at which point it chooses a single \txt{Prepare} message to respond to, ensuring that they are participating in a single consensus instance for participant $p$. This is accomplished by the $\msc{Choose}()$ function. However, note that the validator may have responded to multiple \txt{Prepare} messages for different participants in other consensus instances. For the remaining messages, validators respond with a \txt{PrepareNil} message. 

The operator collects responses from validators until a quorum of either \txt{Prepare} or \txt{PrepareNil} messages are obtained, at which point it can propose a Tesseract that extends the latest states of the participants from whose contexts a quorum of successful \txt{Prepare} responses are obtained. 


\subsubsection{Propose}

\cref{alg:step2} describes the \txt{Propose} stage of the protocol. Among the responses to the \txt{Prepare} messages obtained, the operator identifies the highest states and removes the outdated state information. This is accomplished by the function $\msc{RemoveOutdated}()$. The operator then proposes a Tesseract that extends these states and sends a \txt{Propose} message by forming an ICS by sampling a stochastic set along with the interaction's context. 

In the first view of the very first epoch, the operator may not have any previous state information to extend. In this case, the operator can propose a Tesseract that contains only the interaction itself, without extending any previous states.

Note that the extension is possible if there is no existing Tesseract that has been already agreed upon by the validators in prior views. If such a Tesseract exists, then the operator will re-propose the same Tesseract, ensuring that the consensus instance is consistent with the previous state. 
Further, the proposal is made only for Tesseracts with the same participant set from which the latest state information was obtained.

Validators accept the proposal only if it is \emph{valid}, resulting from the extension of states after removing outdated ones. %They respond to the operator with a \txt{Vote} message.
%Once the operator receives a \emph{supermajority} of \txt{Propose} messages, it forms a \emph{quorum certificate} (QC) and sends a \txt{Vote} message to the interaction's context. This QC represents the operator's commitment to the proposed Tesseract.

% In this section, we describe our protocol Krama in detail. 
% Krama operates in one of two modes: normal mode within each epoch and a reconfiguration mode between epoch transitions. 
% % We first detail the normal operation of our protocol and then the epoch reconfiguration procedure. 
% In what follows, we first establish some basic terminology and proceed to describe the protocol during normal operation and we conclude by detailing epoch reconfiguration. %describe the two modes in detail after establishing some basic terminology.
\begin{breakablealgorithm}
\caption{The Protocol: Stage 2 - Propose}\label{alg:step2}
\begin{algorithmic}[1]         
\LComment{\textbf{\color{black}As an operator:}}
\For{quorum of \txt{Prepare} or \txt{PrepareNil} messages from $\msc{ctx}(\m{ixn}_v)$}
    \State $P\gets\{p\mid \mathrm{received\ supermajority\ of\ \txt{Prepare}\ messages\ from\ \calC_p}\}$
    \State $M\gets$ quorum of \txt{Prepare} messages from $\bigcup\limits_{p\in P}\calC_p$
    % \ForAll{$p\in P_{\m{ixn}_v}$}
    %     \If{received supermajority of \txt{Prepare} messages $M'$ from $\calC_p$}
    %         \State $Q \gets Q\cup \{p\}$
    %         \State $M\gets M\cup M'$
    %     \EndIf
    % \EndFor
    \State $\m{prepareQC} \gets \msc{QC}(M)$
    \State $\m{qcSet}_v \gets \msc{LockedQCs}(\m{prepareQC})$
     \State \Call{RemoveOutdated}{$\m{qcSet}_v$}
    
    %\If{$\exists$ distinct $\m{qc}_1, \m{qc}_2 \in \m{qcSet}_v : \m{qc}_1.\m{stage} = \m{qc}_2.\m{stage} = \m{prevote}$}
        %\LComment{\textbf{\color{black}Conflicting prevote locks detected, propose nil}}
        %\State $T \gets \m{nil}$
    \If{$\exists$ unique $\m{qc} \in \m{qcSet}_v : \m{qc}.\m{stage} = \m{vote} \land \m{qc}.\m{participants} = P$}
        %\LComment{\textbf{\color{black}Single prevote lock matching current participant set}}
        \State $T \gets \msc{Tesseract}(\m{qc}.\m{ixn})$
    \ElsIf{$\nexists \m{qc} \in \m{qcSet}_v : \m{qc}.\m{stage} = \m{vote}$}
        %\LComment{\textbf{\color{black}No prevote locks, extend from highest committed states}}
        %\State $\m{highQCSet}_v \gets \{\m{qc} \in \m{qcSet}_v \mid \m{qc} = \m{highQC}_p \text{ for some } p \in \msc{participants}(\m{ixn}_v)\}$
        \State $\calI' \gets \msc{ExtendCommitted}(\m{qcSet}, \calG_v, P)$
        \State $T \gets \msc{Tesseract}(\calI')$
    \Else
        \State $T \gets \m{nil}$
    \EndIf
    \State \Call{Propose}{$T, \m{prepareQC}$}
\EndFor
\item[] 

\LComment{\textbf{\color{black}As a validator in the ICS:}}
\For{\msg{\txt{Propose}, T, \m{prepareQC}, \calS, \m{rand}, \m{proof}, \m{view}} from $o_T$}
    \State $\m{qcSet}_v \gets \msc{UniqueLockedQCs}(\m{prepareQC})$
    % \If{$T=\m{nil}$}
        
    % \EndIf
    \If{$ T \neq \m{nil} \land \msc{ValidPropose}(T, \m{qcSet}_v, \calS, \m{rand}, \m{proof}, \m{view})$}
        % \State \Call{UpdateState}{$o_T, T, \msc{ICS}(T, \calS)$}
        \State \textbf{send} \msg{\txt{Vote}, T, \m{view}_v} to $\m{operator}_v$
    \EndIf
    % \State $\msc{UpdateLocks}(\m{prepareQC})$
\EndFor
\item[]

\For{\msg{\txt{Propose}, \m{nil}, \m{prepareQC}, \m{view}} from $\m{op}$}
    % \State $\msc{UpdateLocks}(\m{prepareQC})$
    \State Terminate engaging in consensus for the interaction
\EndFor
\item[]

\LComment{\textbf{\color{black}Procedures:}}
\Procedure{Propose}{$T, \m{prepareQC}$}
    \If{$T=\m{nil}$}
        \State \textbf{multicast} \msg{\txt{Propose}, T, \m{prepareQC}, \m{view}_v} to $\msc{ctx}(T)$   
    \Else
        \State $(\calS_v, \m{rand}_v, \m{proof}_v) \gets$ \Call{SelectStochastic}{$T, \m{view}_v$}
        % \State \Call{UpdateState}{$\m{self}, T, \msc{ICS}(T, \calS_v)$}
        \State \textbf{multicast} \msg{\txt{Propose}, T, \m{prepareQC}, \calS_v, \m{rand}_v, \m{proof}_v, \m{view}_v} to $\msc{ICS}(T, \calS_v)$   
    \EndIf
\EndProcedure
\item[] 

%\Procedure{UpdateState}{$\m{op}, T, \m{ICS}$}
%   \State $\m{operator}_v \gets \m{op}$
%   \State $\m{tesseract}_v \gets T$
%    \State $\m{ICS}_v \gets \m{ICS}$
%\EndProcedure
%\item[]

% \Procedure{RemoveOutdatedQCs}{$\m{qcSet}$}
%     \State $\msc{ComputeHighQCs}(\m{qcSet})$
%     %\State $\{\m{maxHeightPrecommit}_p\mid \exists \m{qc}\in \m{qcSet} : p\in \m{qc}.\m{participants}\} \gets \msc{MaxHeight}(\m{qcSet}_v, \m{precommit})$
%     %\State $\m{maxHeightPrevote}_v \gets \msc{MaxHeight}(\m{qcSet}_v, \m{prevote})$
%     \ForAll{$\m{qc} \in \m{qcSet}$ such that $\m{qc}.\m{stage} = \m{vote}$}
%         \ForAll{$p \in \m{qc}.\m{participants}$}
%             \If{$\m{qc}.\m{height}(p) \leq \m{highCommitQC}_p.\m{height}(p) \lor \m{qc}.\m{height}(p) < \m{highVoteQC}_p.\m{height}(p)$}
%                 \State $\m{qcSet} \gets \m{qcSet} \setminus \{\m{qc}\}$
%                 \State \textbf{break}
%             \ElsIf{$\m{qc}.\m{height}(p) = \m{highVoteQC}_p.\m{height}(p) \land \m{qc}.\m{view} < \m{highVoteQC}_p.\m{view}$}
%                 \State $\m{qcSet} \gets \m{qcSet} \setminus \{\m{qc}\}$
%                 \State \textbf{break}
%             \EndIf
%         \EndFor
%     \EndFor
% \EndProcedure
%\item[]

% \Procedure{getUniqueLockedQCs}{$\m{prepareQC}$}
%     \State $\m{qcSet} \gets \emptyset$
%     \ForAll{$\m{lockedQC} \in \m{prepareQC}$}
%         \If{$\m{lockedQC}.\m{ixn} \notin \{\m{qc}.\m{ixn} \mid \m{qc} \in \m{qcSet}\}$}
%             \State $\m{qcSet} \gets \m{qcSet} \cup \{\m{lockedQC}\}$
%         \EndIf
%     \EndFor
%     \State \Return $\m{qcSet}$
% \EndProcedure
% \item[]
% \Procedure{getMaxHeight}{$\m{qcSet}, \m{stage}$}
%     \State $\m{maxHeight} \gets$ map initializing all heights to $-\infty$
%     \ForAll{$\m{qc} \in \m{qcSet}$ such that $\m{qc}.\m{stage} = \m{stage}$}
%         \For{all $p \in \m{qc}.\m{participants}$}
%             \State $\m{maxHeight}_p \gets \max(\m{maxHeight}_p, \m{qc}.\m{height}(p))$
%         \EndFor
%     \EndFor
%     \State \Return $\m{maxHeight}$
% \EndProcedure
\end{algorithmic}
\end{breakablealgorithm}

\begin{algorithm}[hptb]
\caption{The Protocol: Stage 3 - Vote}\label{alg:step3}
\begin{algorithmic}[1]

\LComment{\textbf{\color{black}As the interaction's operator:}}
\For{quorum of \txt{Vote} messages $M$ from $\m{ICS}_v$}
% \msg{\txt{PreVote}, \m{interaction}_v, \m{viewNumber}_v}
    \State $\m{voteQC} \gets $ \sc{QC}($M$)
    % \State Aggregate the quorum of votes into $\m{prevoteQC}$
    \State \textbf{multicast} \msg{\txt{Vote}, \m{tesseract}_v, \m{view}_v, \m{voteQC}} to $\m{ICS}_v$
\EndFor
\item[]

\LComment{\textbf{\color{black}As a validator in the ICS:}}
\For{\msg{\txt{Vote}, \m{tesseract}_v, \m{view}_v, \m{voteQC}} from $\m{operator}_v$}
    \ForAll{$p\in P_{\m{tesseract}_v}$ such that $\m{self}\in \calC_p$}
        \State $\m{lockedQC}_v^p \gets \m{voteQC}$
    \EndFor
    \State \textbf{send} \msg{\txt{Commit}, \m{tesseract}_v, \m{view}_v} to $\m{operator}_v$
    %\EndIf
\EndFor


\end{algorithmic}
\end{algorithm}



\begin{algorithm}[hpbt]
\caption{The Protocol: Stage 4 - Commit}\label{alg:step4}
\begin{algorithmic}[1]

\LComment{\textbf{\color{black}As the interaction's operator:}}
\For{quorum of \txt{Commit} messages $M$ from $\m{ICS}_v$}
% \msg{\txt{PreCommit}, \m{interaction}_v, \m{viewNumber}_v}
    \State $\m{commitQC} \gets $\Call{QC}{$M$}
    % \State Aggregate the quorum of votes into $\m{precommitQC}$
    \State \textbf{multicast} \msg{\txt{Commit}, \m{tesseract}_v, \m{view}_v, \m{commitQC}} to $\m{ICS}_v$
    %\State Execute $\calI_P$ and commit the resulting tesseract
    \State \Call{commit}{$\m{tesseract}_v$}
    %\State Terminate engaging in consensus for $\calI_P$
\EndFor
\item[]

\LComment{\textbf{\color{black}As a validator in the ICS:}}
\For{\msg[]{\txt{Commit}, \m{tesseract}_v, \m{view}_v, \m{commitQC}}}
    \State \Call{commit}{$\m{tesseract}_v$}
    \ForAll{$p\in P_{\m{tesseract}_v}$ such that $\m{self}\in \calC_p$}
        \State $\m{lockedQC}_v^p \gets \m{commitQC}$
    \EndFor        
    %\State Terminate engaging in consensus for $\calI_P$
\EndFor
\item[] 

\For{$\m{view}.\m{end}$}
    \State Terminate engaging in consensus for all interactions
\EndFor

\end{algorithmic}
\end{algorithm}


\subsubsection{Vote and Commit}


The next two stages of the protocol are \txt{Vote} and \txt{Commit} phases, as described in \cref{alg:step3} and \cref{alg:step4} respectively. Validators that accept a proposal respond with a \txt{Vote} message and refrain from voting on other tesseracts with the same participant set. Once a quorum of validators accept the proposal, the operator forms a \emph{quorum certificate} (QC) of the \txt{Vote} responses and broadcasts it to the ICS. Validators on receiving this \emph{lock} onto the Tesseract and pass it in any subsequent \txt{Prepare} phases. This lock represents the validator's commitment to the proposed Tesseract. 
They further respond to the operator with a \txt{Commit} message and the operator, upon receiving a quorum of such messages, forms a quorum certificate and sends a \txt{Commit} message to the ICS as well as the entire validator set. 
Validators then commit the Tesseract to their histories and lock onto the QC. 

\subsection{Reconfiguration mode}

At the beginning of each epoch, validators enter a \emph{reconfiguration mode} to update the validator list and prepare for the next epoch. This reconfiguration process is crucial for maintaining the integrity and security of the protocol, as it allows validators to join or leave the network and ensures that the validator set remains up-to-date.

Two steps outline the reconfiguration process: validator list updation and context shuffling. The former updates the list of validators for the next epoch, while the latter ensures that the contexts of the validators are shuffled to prevent collusion and bribery attacks.

\subsubsection{Validator list updation}

During the epoch, validators may signal to the network their intent to leave or join. The requests are registered and updated in the local state at the beginning of the next epoch. Additionally, the ordering between validators are updated based on the latest validator list.

\subsubsection{Context shuffling}

The context of each participant will remain fixed for the duration of the epoch. However, at the beginning of each epoch, the context of each validator is shuffled to ensure that validators cannot collude with each other or engage in bribery attacks.
The shuffling procedure follows the swap-or-not shuffle described in \cite{HMR12}, which is also detailed in the Ethereum 2.0 specifications and technical handbook \cite{Edg25}. This shuffle allows for the computation of the shuffled position of a single index in a list without requiring the entire list to be shuffled, making it efficient for the protocol's needs.
% \begin{algorithm}
\caption{Context Shuffling}\label{alg:shuffle}
\begin{algorithmic}[1]

% \For{$\m{epoch}.\m{end}$}
%     \State Remove validators that exited in the previous epoch
%     \State Assign new validators that joined over the previous epoch to available indices and new indices at the end of the list
% \EndFor
% \item[] 

\Procedure{ShuffleContext}{$p, \m{seed}, \bbV$}
    \State $\calC_{p}^{'} \gets \emptyset$
    \ForAll{$\m{v} \in \calC_p$}
        \State $\m{index} \gets \bbV.\m{index}(v)$
        \State $\m{shuffledValidator} \gets\ $\Call{ShuffleIndex}{$\m{index}, \bbV$} 
        \State $\calC_{p}^{'} \gets  \calC_{p}^{'} \cup \{\m{shuffledValidator}\}$
    \EndFor 
    \State $\calC_{p} \gets \calC_{p}^{'}$
\EndProcedure
\item[] 

\Procedure{ShuffleIndex}{$\m{index}, \bbV$}
    \State $n \gets |\bbV|$
    \ForAll{$i \in \{1,2,\ldots, \m{numRounds}\}$} %\Comment{$\m{numRounds}$ is a protocol parameter}
        \State $\m{roundSeed} \gets \m{hash}(\m{seed}||i)$
        \State $\m{pivot} \gets \m{roundSeed} \mod n$
        \State $\m{reflection} \gets (\m{pivot} -\m{index} + n)\mod n$
        \If{$\m{hash}(roundSeed||index) = 1$}
            \State $\m{swap}(\m{validator}(\m{index}), \m{validator}(\m{reflection}))$
        \EndIf
    \EndFor
    \State \Return $\m{validator}(\m{index})$
\EndProcedure


% \Procedure{SwapOrNotShuffle}{$ \m{roundSeed}, \m{list}$} 
%     \State $\m{n} \gets \m{len}(list)$
%     \State $\m{pivot} \gets \m{hash}(\m{roundSeed}) \mod \m{n}$
%     \State $\m{mirror} \gets \m{pivot}/2$
%     \State $\m{count} \gets 0$
%     \ForAll{$i \in (\m{mirror}, \m{pivot})$}
%         \State $\m{count} \gets \m{count}+1$
%         \State \Call{SwapOrNot}{$i, \m{mirror}, \m{roundSeed}, \m{count}$}
%     \EndFor
%     \State $\m{secondMirror} \gets \m{mirror} + n/2$
%     \ForAll{$i \in (\m{pivot}, \m{secondMirror})$}
%         \State $\m{count} \gets \m{count}+1$
%         \State \Call{SwapOrNot}{$i, \m{secondMirror}, \m{roundSeed}, \m{count}$}
%     \EndFor
% \EndProcedure
% \item[]

% \Procedure{SwapOrNot}{$\m{index}, \m{mirror}, \m{roundSeed}, \m{count}$}
%     \If{$\m{index} > \m{mirror}$}
%         \State $\m{reflection} \gets \m{mirror} - (\m{index} - \m{mirror})$
%     \ElsIf{$\m{index} < \m{mirror}$}
%         \State $\m{reflection} \gets \m{mirror} + (\m{mirror} - \m{index})$
%     \EndIf
%     \State $\m{coinToss} \gets \m{hash}(roundSeed||count)[1]$
%     \If{$\m{coinToss} \gets 1$}
%         \State $\m{swap}(\m{validator}(\m{index}), \m{validator}(\m{reflection}))$
%     \EndIf
% \EndProcedure

\end{algorithmic}
\end{algorithm}
% This section describes the reconfiguration procedure during epoch transitions. 
% This procedure is initiated towards the end of each epoch. 
% %Validators joining after the start of this procedure are included only at the subsequent epoch's end.


% \subsubsection{Validator List Updation}
% The epoch reconfiguration procedure begins by removing validators who expressed a desire to leave Krama from the validator list. Simultaneously, validators who wished to join are added. This process results in a new list of validators, denoted as $\bbV$, for the upcoming epoch. This list is determined by the validator changes observed during the epoch.

% \subsubsection{Context Shuffling}
% When a participant performs its first interaction in a new epoch, the interaction's operator modifies it to include an updated context as part of the state transition. 
% An ICS formed using the interacting participants' contexts from the previous epoch is used to agree upon the resulting tesseract.
% %The resulting tesseract is agreed upon by an ICS formed using the interacting participants' contexts from the previous epoch. 
% This updated context is then employed for the remainder of the epoch to agree on the participant's interactions.
% The procedure done by the operator to compute the new context is given below. 
% The remainder of this section details the procedure undertaken by the operator to compute the new context. 

% The swap-or-not shuffle described in \cite{HMR12} is used to update the context. 
% %The context is updated by the swap-or-not shuffle described in \cite{HMR12} originally in the context of building enciphering schemes. 
% A description of the same is also present in the Ethereum 2.0 specifications and its technical handbook \cite{Edg25}. 
% An important property of this shuffle is that it can compute the shuffled position of a single index in a list without requiring to shuffling the entire list.  
% This is particularly advantageous as the protocol shuffles only a single participant's context at any given time, which consists of only a few indices. 
% The shuffling process requires a pseudo-random seed sourced from the distributed randomness beacon as input.%, which is sourced from the distributed randomness beacon. 

% The shuffle occurs over multiple rounds, with each round proceeding as follows. 
% First, a pivot index is randomly selected.
% A mirror index is determined by the pivot along which the reflection of the index to be shuffled is computed.
% A random decision is made whether or not to swap the index and its mirror.
% If swapping, the validators in both indices exchange positions; otherwise, the process moves on to the next round.
% The hash of the input seed concatenated with the round number serves as the randomness required for selecting the pivot. 
% The randomness required for selecting the pivot each round is given by the hash of the input seed by concatenated with the round number. 
% The decision to swap or not is determined by the hash of the input seed concatenated with both the round number and the index to be shuffled. 
% The pseudocode for the above procedure is given in \Cref{alg:shuffle}.

% \subsubsection{Context Reset}

% Participants in Krama can initiate a special interaction requesting their context to be monitored due to inactivity of context validators.
% %In Krama, participants have the option to initiate a special interaction requesting monitoring of their context due to inactivity of context validators. 
% This process unfolds as follows. 

% Validators, upon receiving this type of interaction, perform an additional check to determine if the time elapsed since the previous tesseract formation indicates a sufficiently high probability of context validator unavailability.
% This interaction is agreed upon by an ICS with sarga's context. 
% Validators, upon receiving the resulting tesseract monitor the liveness of the participant's context validators at regular intervals for a specified period. 
% During this monitoring period, if the context validators are found unavailable, validators broadcast a vote for resetting the participant's context.
% Otherwise, they vote against it.
% A separate context reset logic present with each validator manages this process.
% The participant can query the context reset logic to determine if a supermajority of votes favour resetting their context.
% On receiving a supermajority, the participant can broadcast a \emph{context reset} interaction. 
% This interaction includes sarga and similar to the regular context shuffling procedure during epoch reconfiguration, the participant's context is shuffled, agreed upon, and the resulting tesseract is committed.

% \subsubsection{Validator Weights}
% In the actual implementation of Krama, validators have weights that determine their voting power to prevent Sybil attacks, see \cite{PS17b}. 
% Thus, the protocol slightly differs from the above, with the definition of a supermajority, and the stochastic set sampling and context shuffling procedures using the validators' weights.
% %Thus, a supermajority over a set of validators with total weight $w=3b+1$ is a subset of validators with total weight $2b+1$.
% %Further, the protocol slightly differs from above with the stochastic set sampling and context shuffling procedures using the validators' weights.
% This gives validators with more weight a higher chance to participate in consensus.

\subsection{Utility methods}

\subsubsection{Operator selection}

Deterministic selection of the operator for an interaction is crucial for ensuring that all validators agree on the same operator. This selection is based on the participant set of the interaction and the current view. The operator is chosen from the context of the participants in a round-robin manner, ensuring that each participant has a chance to drive consensus in different views.

% Each participant is designated a publicly known validator called an \emph{operator} to drive consensus in each view. 
% Operators for a view are selected from the participant's context validators using round-robin assignment to views within an epoch. 
% This mapping is established during reconfiguration at the epoch's beginning.
% We denote the operator for participant $p$ in view $v$ as $\msc{operator}(p,v)$.
% For an interaction $I$ in view $v$, the set of operators for the interacting participants is $\mathcal{O}_I=\{\msc{operator}(p,v)\mid p\in P_I\}$.
% A deterministic procedure selects a single operator from $\mathcal{O}_I$ to drive consensus for interaction $I$, denoted by $\msc{operator}(P_I, v)$.

\subsubsection{Stochastic set of validators}

\begin{algorithm}[hptb]
\caption{Sampling the Stochastic Set}\label{alg:sample}
\begin{algorithmic}[1]         

\Procedure{SelectStochastic}{$T, \m{view}$}
    %\State $P \gets \msc{participants}(\calT)$
    \State $\mathit{\bbV'} \gets \bbV\setminus\bigcup\limits_{p\in P_T}\calC_p$
    \State $\m{seed} \gets \mathrm{XOR}(\,\m{hash}(T), \msc{random}(\m{view})\,)$
    \State $(\m{rand}, \m{proof}) \gets \mathit{VRF}(\mathit{seed}, \mathit{sk}_\self)$ 
    \State $\calS \gets \emptyset$
    \ForAll{$i \in \{1,2,\ldots,n_\calS\}$}
        \State $\m{index} \gets \m{hash}(\m{rand}||i)$ \Comment{$\m{index} \in [1, n_V]$}
        \State $\calS \gets \calS \cup \{\m{\bbV'}[\m{index}]\}$
    \EndFor
    \State \Return $(\calS, \m{rand}, \m{proof})$
\EndProcedure

\end{algorithmic}
\end{algorithm}


To decide on an interaction, the operator needs to form an ICS. A key component is the stochastic set of validators. The operator samples uniformly according to~\cref{alg:sample} after excluding the validators part of a participant's context. The sample is associated with a proof obtained from a verifiable random function (VRF) to ensure that the sampling is verifiable by other validators.


% is sampled uniformly at random from all validators excluding those in the interacting participants' contexts. 
% This random sampling minimizes bribery and collusion attacks by requiring attackers to corrupt a significant proportion of the entire validator set.
% %Further, the condition requiring a quorum to consist of a supermajority of votes from each set in the ICS including the stochastic set ensures that such attacks made only towards the participants' contexts do not affect safety.

% The stochastic set is sampled as follows. 
% First, eligible validators are computed by excluding those in the interacting participants' contexts. 
% A seed is generated by XORing the tesseract hash with the view's random number from the distributed randomness beacon.
% The operator runs a VRF with this seed, producing a pseudo-random output and proof. 
% The stochastic set of $n_\calC$ validators is selected from eligible validators using $n_\calC$ random numbers generated by hashing the pseudo-random output concatenated with sequential digits. 

% Validators can verify the stochastic set by repeating the operator's steps using the pseudo-random output. 
% The VRF's verifiability property allows verification of the pseudo-random output using the proof and operator's public key.

\subsubsection{Accepting the proposal for an interaction}\label{validpropose}

Each validator that is part of a \emph{consensus instance} receives a proposal from the operator and has to decide whether to accept or reject the proposal. 
The decision is based on the following criteria:

\begin{enumerate}
    \item The validator belongs to the corresponding ICS formed with the interacting participants' contexts and the stochastic set. 
    %\item The quorum certificate $\m{prepareQC}$ is valid.
    \item The sender is the valid operator for the tesseract. 
    \item The validator did not \txt{Vote} on a different tesseract with any of the interacting participants during the current view. 
    \item The proposed tesseract is the valid output of the state extension procedure undertaken by the operator. 
    %\item The proposed tesseract extends the validator's state.
    \item The proposed stochastic set contains validators sampled according to the stochastic set selection procedure. 
\end{enumerate}

The logic for checking the criteria is abstracted in the function $\msc{ValidPropose}()$ in the pseudocode.

% \subsubsection{Initialisation of New Participants}

% A newly joined participant's context is initialised through a special interaction agreed upon by sarga's operator, who forms an ICS with sarga's context and a stochastic set. 
% This stochastic set becomes the new participant's context. 
% Upon agreement, the tesseract is committed by linking it to sarga's latest state with edges for both participants.

% The operator undertakes the following procedure to determine a proposal to participant set $P$ that maintains safety and liveness. 
% First, the operator obtains all the $\m{lockedQCs}$ present in $\m{prepareQC}$ corresponding to prevote and precommit stage locks on tesseracts, and stores them in a set $\m{qcSet}$. 
% This is done by the function $\msc{LockedQCs}(\m{prepareQC})$ in the pseudocode.
% %For each participant $p$ involved in an interaction within $\m{qcSet}$, the prevote and precommit stage QCs of the highest height and latest view are identified.
% The operator identifies the prevote and precommit stage QCs of the highest height and latest view for each participant $p$ involved in an interaction within $\m{qcSet}$.
% %This is denoted as $maxHeightPrevote_p$ and $maxHeightPrecommit_p$ for a participant $p$.
% This is denoted as $highPrevoteQC_p$ and $highPrecommitQC_p$ for a participant $p$. 
% Now, the operator removes prevote x	 locks from $\m{qcSet}$ that are superseded by committed tesseracts or locks from later views. 
% This is ascertained for a prevote stage $\m{lockedQC}$ by comparing the height of each participant in the corresponding tesseract with the heights of the previously identified highest QCs as well as the view number of the highest prevote QCs. 
% This is done as follows for a tesseract $T$ corresponding to a prevote stage $\m{lockedQC}$.
% If there exists a participant $p\in P_T$ such that $T.\m{height}(p)\leq \m{highPrecommitQC}_p.\m{height}(p)$ or $T.\m{height}(p)<\m{highPrevoteQC}_p.\m{height}(p)$, then a committed tesseract is present at the same height for participant $p$ and thus, the tesseract $T$ is removed from $\m{qcSet}$. 
% Here, for a quorum certificate $\m{qc}$, $\m{qc}.\m{height}(p)$ refers to the height for the tesseract corresponding to the quorum certificate. 
% Also, as a liveness rule, if the tesseract $T$ has a participant $p\in P_T$ such that the prevote stage QC for $T$ was formed at an earlier view than $highPrevoteQC_p$, then the prevote stage QC is removed from $\m{qcSet}$.
% Thus, we safely ignore such locks and remove them from $\m{qcSet}$
% It is crucial for liveness to perform this procedure. 
% The procedure $\msc{RemoveOutdatedQCs}()$ in the pseudocode encapsulates this process. 
% Based on the remaining QCs in $qcSet$, the operator selects an interaction to propose according to the following rules:
% \begin{enumerate}
%     \item If exactly one prevote stage QC remains in $qcSet$ and it has participant set $P$, the operator proposes the interaction associated with this QC.
%     \item If no prevote stage QCs remain in $qcSet$, the operator proposes a new interaction that is a $p$-child of the tesseract with the current state of each participant $p\in P$. This involves taking the highest $\m{precommitQC}$/committed tesseract for each participant and proposing an interaction that extends these states. This is done by the function $\msc{ExtendCommitted}()$.
%     \item If neither of the above conditions are met, the operator refrains from proposing any interactions and instead sends the propose message with the $\m{nil}$ tesseract as given in more detail below.
% \end{enumerate}
% The operator cannot propose an interaction with a different participant set due to its lack of knowledge about the latest lock information of other participants and the participant context being on a different consensus instance. 
% If the above procedure results in a valid interaction, the operator executes it and proposes the resulting tesseract. 
% The proposal is broadcast to an ICS formed by sampling a stochastic set as described in \Cref{subsubsec:randsetsample}. 
% This \txt{Propose} message contains (1) the tesseract, (2) the quorum certificate $\m{prepareQC}$, (3) the randomly selected stochastic set, and (4) the corresponding pseudo-random number and its proof. 
% However, if the above procedure does not result in an interaction for proposal, the operator broadcasts a \txt{Propose} message with the $\m{nil}$ tesseract and quorum certificate $\m{prepareQC}$ to the interaction’s context.
% This message serves to share the $\m{prepareQC}$ with the interaction’s context
% so that the validators are aware of the latest locked QCs of the interacting participants. 
% The operator can only propose a single tesseract with a given participant involved during a view, preventing conflicting interaction proposals.

% A validator, on receiving a \txt{Propose} message with a valid tesseract, verifies whether the following conditions ($\dagger$) are satisfied.
% \begin{enumerate}
%     \item The validator belongs to the corresponding ICS formed with the interacting participants' contexts and the stochastic set.
%     %\item The quorum certificate $\m{prepareQC}$ is valid.
%     \item The sender is the operator for the tesseract during the current view.
%     \item The validator did not \txt{PreVote} on a different tesseract with any of the interacting participants during the current view. 
%     \item The proposed tesseract is the valid output of the procedure undertaken by the operator.
%     %\item The proposed tesseract extends the validator's state.
%     \item The proposed stochastic set contains validators sampled according to the stochastic set selection procedure.
% \end{enumerate}
% The $\msc{validPropose}$ predicate in \Cref{alg:step2} evaluates to true if the above conditions are satisfied. 
% On satisfying the above, the validator responds to the operator with a \txt{PreVote} message. 
% Now, during the current view, the validator cannot engage in consensus (i.e. send \txt{PreVote} or \txt{PreCommit} messages) for any other tesseracts whose participant set overlaps with the current interaction's participant set.
% This restriction maintains the following invariant: A validator votes (\txt{PreVote} and \txt{PreCommit}) for at most one tesseract with a given participant in each view. 
% %cannot vote (\txt{PreVote} and \txt{PreCommit}) for two different interactions with a common interacting participant in a given view. 
% However the validator refrains from voting if the \txt{Propose} message contains the $\m{nil}$ tesseract.

% The validator concludes by updating its locks using $\m{prepareQC}$. 
% If any tesseracts corresponding to precommit stage QCs are missing, the validator retrieves them from other validators in the network.
% For each participant $p$, if the validator is present in its context and the current $\m{lockedQC}_p$ is from the prevote stage and a prevote stage QC from a higher view is present in $\m{prepareQC}$, then the validator updates $\m{lockedQC}_p$ to the higher view QC.
% The function $\msc{UpdateLocks}(\m{prepareQC})$ performs this updation.


% -----
% Old text below this.

% \subsection{The Basics}

% \subsubsection{Validator List}
% There is a total order on the set of validators $\bbV$ and is therefore maintained as a list.
% This list is updated during epoch transitions as part of the reconfiguration procedure. % by shuffling the positions of the validators in the order. 
%During epoch transitions, this set and the total order is updated
%As defined in the preliminaries, $\bbV$ is the set of validators that participate in consensus.  \srinidhi{Is it important to mention that it is a list?}
%Between epoch transitions, the set may receive updates as part of the reconfiguration procedure.
% list undergoes updates during each epoch transition as part of the reconfiguration procedure. 

% \subsubsection{View}
% The normal mode of Krama proceeds in a sequence of monotonically increasing views where the duration of each view \srinidhi{is this symbolic time like ticks or actual wall clock time?} is $\m{viewLength}$. 
% Validators determine the current view using their local clock, referencing the time of the previous epoch transition. 
% Each validator receives a random number from the distributed randomness beacon at the start of each view.
% Validators attempt to reach consensus on tesseracts within each view. %, and attempt to reach consensus. 
% This consensus process is limited to the view, at the end of which validators stop participating in consensus on tesseracts.
%The consensus process for each tesseract is governed by a timeout called $\m{consensusTimeout}$, which determines how long validators will attempt to reach agreement.
%The duration of this timeout is initially set to $\m{viewLength}$.
%Ideally, validators achieve consensus on a tesseract within the $\m{consensusTimeout}$ duration. 
%The initial timeout of $\m{viewLength}$ duration is set to be sufficient for achieving consensus within a view after GST. 
%Once this timeout is reached, validators stop participating in consensus on a tesseract. 
%\srinidhi{Is this fair to say?} For future interactions involving the same participant set, the consensus timeout duration is increased by one $\m{viewLength}$. 
%This extended duration helps validators combat potential network delays within the interaction's context.
%Views in Krama synchronize the entry of validators into consensus instances. 
% Since multiple such procedures can run simultaneously in Krama, validators need coordinated entry points.
% As views drive this consensus procedure, all validators are available for future participation at the beginning of each view, in sync.
%validators enter new consensus instances in sync, ensuring availability for future participation.
%As consensus duration is always a multiple of $\m{viewLength}$, validators enter new consensus instances in sync, ensuring availability for future participation.

% The consensus duration for a tesseract is determined by a timeout called $\m{consensusTimeout}$, whose duration is initially set to $\m{viewLength}$. 
% If validators are unable to achieve consensus within the stipulated time, they retry after incrementing the timeout by $\m{viewLength}$ \srinidhi{Is this fair to say?} 
% Validators stop participating in consensus on the tesseract once this timeout is triggered. 

% When the consensus timeout is reached, all validators stop participating in consensus for that interaction. Since consensus duration is always a multiple of $\m{viewLength}$, validators enter into achieving consensus on new interactions in sync, ensuring their availability for participating in consensus on future interactions.
% If consensus is not achieved on the proposed tesseract within the allocated time, any future interaction with the same participant set will have its consensus timeout duration increased by one $\m{viewLength}$. This extended duration helps validators combat potential network delays among the interaction's context.

% \subsubsection{Operator}
% A publicly known validator, called an \emph{operator}, is designated for each participant to drive the consensus procedure in a view.
% These operators are obtained from the participant's context. 
% Thus, for each participant, validators from its context are mapped to views of the epoch in a round-robin manner to act as operators.
% Each participant's context validators are assigned to act as operators for each view in an epoch, in a round-robin fashion. 
% This mapping is set at the beginning of the epoch during reconfiguration.
% The operator for a participant $p$ during view $v$ is given by $\msc{operator}(p,v)$.
% Thus, for an interaction $I$ during view $v$, there is a set of operators $\mathcal{O}_I=\{\msc{operator}(p,v)\mid p\in P_I\}$ associated with the interacting participants. 
% A deterministic procedure selects a single operator from the set $\mathcal{O}_I$ to drive the consensus procedure for the interaction $I$. 
% This is denoted as $\msc{operator}(P_I, v)$ in the pseudocode.
%%% Succinct: %%%


%The operator is fixed for a given epoch and changes in a round-robin fashion among the validators of a given participant. 
% This mapping is established during epoch reconfiguration and follows a round-robin pattern among the validators in the participant's context.
%The operator for a given view and participant is publicly known. 

%At the start of each view, a random output is obtained from the distributed randomness beacon, given by $\msc{random}(\m{view})$. 
%The duration for achieving consensus on a tesseract is a multiple of $\m{viewLength}$, at the end of which, the validators terminate engaging in consensus on the tesseract.
%This is given by a timeout called $\m{consensusTimeout}$ whose length is typically $\m{viewLength}$ but is increased by multiples of $\m{viewLength}$ if network delays are observed among validators in the interaction's context in order to combat it.

% \subsubsection{Operator Selection for an Interaction}
% %In each view, interactions can be executed and their resulting tesseracts agreed upon and committed by honest validators. \srinidhi{passive voice}
% We now describe how a single operator is selected to drive the consensus procedure for an interaction $I$ during a view $v$. 
% As given above, for an interaction $I$ in a view $v$, there is an operator associated with each participant $p\in P_I$.
% An operator is selected to drive the consensus procedure on an interaction as follows.
% %This consensus procedure to agree on an interaction $I$ is driven by an operator selected as follows.  
% The operator for an interaction $I$ is determined by its participant set $P_I$ and the current view $v$.
% As mentioned earlier, for each participant $p\in P_I$, there is an operator associated with it for view $v$ given as $\msc{operator}(p,v)$.
% Thus, we have a 

% An operator is deterministically selected from the set of operators associated with the participant set for the view, $\{\msc{operator}(p, v)\mid p\in P_I\}$.
% The function $\msc{operator}(\m{I}, \m{v})$ returns the operator for the interaction in the view. 
% The function can also be invoked with the tesseract resulting from an interaction. 
% \srinidhi{passive voice. Who invokes and how?}

% \subsubsection{Highest QC}
% Each validator in a participant's context stores the quorum certificate for the highest stage of the highest interaction involving the participant as $\m{highestQC}$.
% This $\m{highestQC}$ is different for each participant as each participant has its own chain.

% During the Prepare phase, validators respond with their $\m{highestQC}$ which gives the highest QC that they have seen for their associated participant. 
% The operator, on receiving a quorum of responses, computes the $\m{highestQC}_p$ for each of the interacting participants as the highest QC found among the responses. 
% The operator then has to propose a block that agrees with the highest QCs for each of the interacting participants. 

%\srinidhi{The whole subsection needs a story, right now it reads like a set of disjoint and disconnected sectences all grouped together. Try starting with a topic sentence to set the context for the paragraph and follow from there. Is this meant to just be a set of definitions? You need to motivate each of them.}


% \subsection{Normal Operation}

% \subsubsection{The Upshot}
% We first present an overview of Krama's normal operation. 
% Validators reach consensus on individual tesseracts containing updated state from executing interactions. 
% In a single view, validators can reach consensus on multiple tesseracts with mutually disjoint participant sets. 
% At the start of each view, operators start a \emph{consensus instance} by sending a message to an interaction's context. 
% This consensus instance ideally leads to a decision with the resulting tesseract committed. 
% This is in contrast to traditional blockchain protocols where several interactions are batched into a single block and agreed upon one at a time in a single global consensus instance.
% Moreover, Krama generates tesseracts only when interactions occur between participants. 
% Normal operation has four stages: Prepare, Propose, Prevote, and Precommit. 
%The Prepare stage is partly governed by a timeout whereas validator actions in the remaining stages are triggered by message reception as well as view timeouts.
% A timeout partly governs the Prepare stage, while validator actions in the remaining stages are triggered by message reception and view timeouts.
%Validator actions are triggered by message reception and view timeouts.
% All validators know the interaction's operator for a given view.
% In the consensus procedure for an interaction, communication occurs only between the designated operator and the interaction's context/ICS. 
% Unlike PBFT-like protocols, Krama uses predetermined views with leaders changing for each new tesseract. 
% Since tesseracts aren't continuously generated, validators cannot predetermine the next consensus-driving operator cannot be ---it exists only when participants send interactions to validators. 
% Thus, each operator requires a communication round-trip with the participants' contexts to propose a safe tesseract extending the interacting participants' latest state.
% The Prepare stage handles this: the operator learns each participant's latest state from the interaction's context and proposes a tesseract extending it. 
% Further, validators in a participant's context receive \txt{Prepare} messages from operators until a timeout is reached. 
% A participant's context can receive multiple such \txt{Prepare} messages and upon reaching the timeout, the validators choose a single \txt{Prepare} message to respond to with the latest state information.
% This ensures that among a set of Prepare messages with overlapping participant sets, at least one consensus instance can successfully reach a decision during the view. 
% The operator makes the proposal to the participants from whose contexts a quorum of successful responses was received. 
% The operator forms an ICS by randomly sampling a stochastic set using VRFs to make this proposal. 
%On receiving a quorum of responses, the operator makes the proposal to an ICS formed by randomly sampling a stochastic set using VRFs. 
%The view's random number is used as the seed for the VRF. 
% Proposals include justification in the form of the latest states obtained.
% Validators respond to valid proposals with \txt{PreVote} messages. 
% Upon receiving a quorum, the operator forms a QC and sends a \txt{PreVote} message with the QC back to the ICS. 
% Validators then lock onto the tesseract and respond with \txt{PreCommit} messages. 
% They persistently store this QC representing the lock for each participant context they belong to. %, sending locked QCs to operators during prepare phases.
% Validators persistently store this lock by storing the corresponding QC for each participant whose context they belong to. 
% This locked QC is sent to the operator during the prepare phase so that the operator is aware of the latest locks. 
% Now, operators on receiving a quorum of \txt{PreCommit} messages form a QC and respond with a \txt{PreCommit} message containing the QC. 
% Validators then commit the tesseract to their histories and lock onto the QC. 
% If the view ends, validators terminate participating in consensus for any tesseracts.
% To ensure liveness, operators include locked QCs from the Prepare phase in \txt{Propose} messages.
% Validators update their locks when encountering higher-view locks in these QCs. 
% If the operator cannot find a tesseract extending the obtained locks during prepare, it sends a \txt{Propose} message without a tesseract but containing the locked QCs for validators to update their locks.




% We now describe the protocol under normal operation in detail. 

% Algorithms 1, 2, 3, and 4 delineate the entire protocol during normal operation, covering each of the four stages. 
% The protocol description is divided according to stages. 
% The protocol's specification is divided into two parts for validators and operators. 
% In the protocol description, signature and quorum certificate verification is implicit on message reception.
% Further, validators only accept messages from the current view.
% %The protocol described below focuses on achieving consensus on a single tesseract.
% %Multiple instances of the below procedure can run simultaneously for different tesseracts with disjoint participant sets. 
% The protocol below focuses on a single tesseract, but multiple instances of the  procedure can run simultaneously for different tesseracts with disjoint participant sets.


% \subsubsection{Notation}
% Before proceeding, we introduce some common notation used in the pseudocode. 
% In the pseudocode, $\m{self}$ refers to the validator performing the operations. 
% Validators maintain several state variables for consensus on a tesseract including: the tesseract itself, the current operator, and the ICS, denoted with a subscript $v$. 
% $\msc{ICS}(T, \calS)$ represents the ICS for a tesseract $T$ with a stochastic set $\calS$.
% The function $\msc{tesseract}()$ returns the tesseract from executing an interaction.
% The function $\msc{QC}()$ returns a \emph{quorum certificate} from a quorum of votes.
% $\msc{view}(T)$ gives the view in which a tesseract $T$ is committed. \todo{This notation is only used in the safety pf and not in the algo.}

% \subsubsection{Prepare Stage}

% %%%OLD: Till 29th May 2025
% The prepare stage serves as the initial phase of each view, functioning as a discovery mechanism. 
% This stage is conceptually similar to the view-change phase in protocols such as PBFT \cite{CL99} and Fast-HotStuff \cite{JNFG23}.
% However, Krama's approach differs in several key aspects:
% \todo[inline]{Need to introduce the specifics of this earlier in the protocol section}
% (1) Views in Krama are predetermined, unlike PBFT-like protocols where view changes occur only when the current leader fails to make progress. 
% (2) As consensus is achieved for each tesseract individually, there is no guarantee that a tesseract shall be extended in the future; Thus, the next consensus-driving operator cannot be predetermined as such an operator shall exist only when a corresponding interaction is sent to the validators by the participants. 
% %and therefore, a round of communication must be dedicated for an interaction's operator to learn about the latest state of the interacting participants each time. 
% %Due to the presence of multiple participant chains and dynamic consensus based on interacting participants, the next consensus-driving operator cannot be predetermined and hence cannot learn about the current consensus state of the participants. 
% Consequently, a round of communication must be dedicated for the operator to acquire the latest state of the participants a each view's commencement.
% This is achieved in the Prepare stage of the protocol. 
% %Hence, the operator obtains knowledge about previously locked tesseracts during the prepare stage at the start of the view.
% %The above passage was rewritten using Notion AI to make it more readable. Date: 18th Oct 2024

% %The pseudocode for this stage is given in \Cref{alg:step1}.

% \subsubsection{As an operator for the view}
% At the start of a view, operators for the view select an interaction ($\m{ixn}$) from the pool of interactions for which they are the operator, given by $\msc{operator}(P_\m{ixn}, \m{view})$. 
% When an interaction corresponding to a prevote stage $\m{lockedQC}$ is present in the pool of interactions, it takes priority on being chosen from the pool.
% This aids in ensuring liveness. \srinidhi{takes priority over what?}
% %The operator proceeds to send a \txt{Prepare} message containing the interaction to the interaction's context.
% The operator proceeds to send a \txt{Prepare} message to the interaction's context.

% \subsubsection{As a validator in the interaction's context}
% Upon receiving a \txt{Prepare} message for an interaction whose context they belong to, validators respond with a \txt{Prepare} message containing their $\m{lockedQC}^p$ from previous views. Here, (1) $p$ is an interacting participant whose context they belong to, and (2) the quorum certificate $\m{lockedQC}^p$ is either a prevote or precommit stage QC, corresponding to the highest view the validator has obtained for participant $p$.
% The validator is essentially \emph{locked} on to the tesseract corresponding to the $\m{lockedQC}$.
% Validators persistently store $\m{lockedQC}^p$ for each participant $p$ whose context they belong to.
% %The above passage was rewritten using Notion AI to make it more readable. Date: 18th Oct 2024
% % Validators, upon receiving a \txt{Prepare} message for an interaction whose context they belong to, respond with a \txt{Prepare} message containing its $\m{lockedQC}^p$ from previous views. 
% % Here, $p$ is the participant in whose context the validator is in where $p$ is an interacting participant for the interaction. 
% % The quorum certificate $\m{lockedQC}^p$ is the prevote or precommit stage quorum certificate of the highest view the validator has obtained for the participant $p$.
% % The validator is essentially \emph{locked} on to the tesseract corresponding to the $\m{lockedQC}$.
% % Validators persistently store $\m{lockedQC}^p$ for each participant $p$ whose context it belongs to.
%%%%
%%%% NEW: 30th May 2025


%On receiving a \txt{Prepare} message for an interaction in whose context they are present, validators respond with the $\m{lockedQC}$ for each interacting participant in whose context they are present. 
%On receiving a \txt{Prepare} message in whose interaction's context the validator is present, a response with the $\m{lockedQC}$ for each interacting participant in whose participant context they are present is sent back to the operator. 
% The prepare stage serves as the initial discovery phase at the start of each view, with context validators deciding on a single operator's ICS to be a part of and sharing their highest locks with the operator.
% Operators for an interaction send \txt{Prepare} messages to the interaction's context. 
% Validators in participants' contexts can receive multiple \txt{Prepare} messages from different operators, each corresponding to a different participant set. 
% For instance, interactions $I_1$ and $I_2$ with participant sets $P_1=\{p_1,p_2\}$ and $P_2=\{p_2, p_3\}$ can have their corresponding operators for a given view from $\calC_{p_1}$ and $\calC_{p_3}$ respectively.
% Here, the context validators of $p_2$ receive \txt{Prepare} messages for both the participant sets $P_1$ and $P_2$. 
% Only one of $I_1$ or $I_2$ can be decided successfully in the current view while maintaining safety.
% Therefore, the context validators of $p_2$ must choose a single consensus instance to participate in to avoid conflicts.
% Thus, among multiple \txt{Prepare} messages obtained, context validators pick and respond to a single chosen one. 
% For this, the stage is governed by a timeout starting at the beginning of the view. 
% Until this timeout, context validators collect \txt{Prepare} messages from operators.
% Context validators choose a single \txt{Prepare} message to respond to when this timeout is triggered.
% Validators order the \txt{Prepare} messages deterministically and choose the highest one. 
%This is given by the \textsc{Sort}() function which sorts the \txt{Prepare} messages obtained, with the highest one being chosen. 
% This results in all context validators with the same set of \txt{Prepare} messages deciding on the same message to respond to. 
% Further, since the messages are totally ordered, the highest element in the order among a set of such overlapping interactions is chosen by all of its interacting participants' contexts. 
% We capture this more precisely as follows. 
% A set of consensus instances are \emph{overlapping} if in its corresponding family of participant sets, for each participant set there exists another participant set whose intersection is non-empty and neither is a subset of another. 
% Note that there cannot be two identical participant sets in a given view as each participant set has a single operator during a view and hence it would form a single consensus instance. 
% A \emph{maximal} set of overlapping consensus instances is maximal with respect to the inclusion order among all sets of overlapping consensus instances. 
% Validators ordering the \txt{Prepare} messages obtained and choosing the highest one ensures the following. %that among a maximal set of overlapping consensus instances, 
% \begin{lemma}\label{lem:prepareOverlap}
%     In each maximal set of overlapping consensus instances with honest operators, at least one consensus instance successfully obtains a quorum of \txt{Prepare} messages from each of its participants' contexts. 
% \end{lemma}

% The proof of the above is given in \Cref{sec:analysis}. 
% Validators respond to the chosen operator with a \txt{Prepare} message containing their $\m{lockedQC}$ for the interacting participant whose context they belong to.
% For all other operators, validators send a \txt{PrepareNil} message signifying that they chose a different operator. 
% This message contains the set of \txt{Prepare} messages received, using which the operator can verify that the \textsc{Choose}() function resulted in a  different message.
% Further, this choice prioritises interactions with a prevote stage $\m{lockedQC}$
% This ensures liveness as locked interactions eventually get committed as operators rotate.



% \subsubsection{Propose Stage}

% We now elaborate on the Propose stage of the protocol.

% %%%OLD: 29th May 2025
% \subsubsection{As an operator for the view:}
% Upon receiving a quorum of \txt{Prepare} messages from the interaction's context, the operator aggregates them into a quorum certificate $\m{prepareQC}$.
% From this QC, the operator identifies the unique lockedQCs (with unique interactions associated with them) present in prepareQC and stores them in a set $qcSet$. \srinidhi{At this point it is unclear why lockedQC is unique? The way quorum certificates are set needs to be explained earlier, this is different from other family of protocols}
% To determine which interaction to propose such that it maintains safety, the operator follows a systematic process:
% \begin{enumerate}
%     \item The operator identifies the maximum height from precommit stage QCs \srinidhi{We haven't got to the precommit stage yet, so there is no way to know what this means.} \rishal{Shall I postpone the description of this process to after describing the normal case operation completely?} for each participant involved in any interaction within $qcSet$, denoted by $maxHeightPrecommit_p$ for participant $p$. Similarly, the maximum height from prevote stage QCs for each participant is identified and is denoted by $maxHeightPrevote_p$. This is done by the function $\msc{ComputeMaxHeights()}$.
%     \item  The operator then filters outdated prevote QCs by removing any prevote stage QC where:
%     \begin{enumerate}
%         \item A participant $p$ in the QC has the same height or higher in some precommit stage QC by comparing with $maxHeightPrecommit_p$ for the participant, or
%         \item A participant $p$ in the QC has a higher height in some other prevote stage QC by comparing with $maxHeightPrevote_p$.
%     \end{enumerate}
%     This filtering ensures that prevote locks that have been superseded by committed tesseracts are ignored.
%     \item Based on the remaining QCs in $qcSet$, the operator selects an interaction to propose according to the following rules:
%     \begin{enumerate}
%         %\item If multiple distinct prevote stage QCs remain in qcMap, the operator proposes nil (no interaction). This occurs when there are conflicting prevote locks that cannot be resolved, preventing progress on either interaction.
%         \item If exactly one prevote stage QC remains in $qcSet$ and its participant set matches the current interaction's participant set, the operator proposes the interaction associated with this QC.
%         \item If no prevote stage QCs remain in $qcSet$, the operator proposes a new interaction that is a $p$-child of the tesseract with the current state of each participant $p\in P$. This involves taking the highest $\m{precommitQC}$/committed tesseract for each participant and proposing an interaction that extends these states. This is done by the function $\msc{ExtendCommitted}()$.
%         \item If neither of the above conditions are met, the operator refrains from proposing any interactions and instead sends the propose message with the $\m{nil}$ tesseract as given in more detail below. 
%     \end{enumerate}
% \end{enumerate}

% Note that the operator cannot propose an interaction with a different participant set as the latest lock information of other participants are not known to the operator. 

% If an interaction satisfying the above conditions is present, it is executed and the resulting tesseract is agreed upon by the protocol as follows.
% To propose the tesseract, an interaction consensus set is formed by sampling a stochastic set as described in \Cref{subsubsec:randsetsample}.
% The operator broadcasts a \txt{Propose} message to the ICS containing (1) the tesseract, (2) the quorum certificate $\m{prepareQC}$, (3) the current view, (4) the randomly selected stochastic set, and (5) the corresponding pseudo-random number and its proof. 
% However, if such an interaction is not present, then the \txt{Propose} message is broadcast to the interaction's context containing (1) the $\m{nil}$ tesseract, the quorum certificate $\m{prepareQC}$, and (3) the current view.
% This message serves to share the $\m{prepareQC}$ with the interaction's context so that the validators are aware of the highest locked QCs of the interacting participants. 
% The operator is limited to proposing a single tesseract with a given participant involved during a view, preventing conflicting interaction proposals. 
% %This prevents conflicting interactions being proposed by the operator. 
% When agreeing on a tesseract, communication occurs only to and from the operator.


%%%NEW: 30th May 2025
% \noindent\paragraph{As an operator for the view}

% When the operator receives a quorum of \txt{Prepare} or \txt{PrepareNil} responses from the interaction's context, it aggregates the \txt{Prepare} messages into a quorum certificate $\m{prepareQC}$ from participant set $P$. 
% Here, the participant set $P$ is the maximal subset of interacting participants from which a quorum of \txt{Prepare} messages are obtained.
% Note that by \Cref{lem:prepareOverlap}, atleast one consensus instance will have received \txt{Prepare} messages from all of its participants' contexts. 
% However, to allow for potentially more interactions to successfully reach a consensus decision, operators can proceed by proposing to a subset $P$ of the initial participant set for which the \txt{Prepare} message is sent. 


% \paragraph{As a validator in the ICS}\label{subsubsec:invariant}



% \subsubsection{PreVote Stage}

% Upon receiving a quorum of \txt{PreVote} messages from the ICS, the operator aggregates them into a quorum certificate $\m{prevoteQC}$
% and broadcasts a \txt{PreVote} message containing it to the ICS. Validators in the ICS respond with \txt{PreCommit} messages and set their $\m{lockedQC}^p$ as $\m{prevoteQC}$ for interacting participants whose context they belong to.
% This step is crucial for safety.

% \subsubsection{PreCommit Stage} 

% Upon receiving a quorum of \txt{PreCommit} messages from the ICS, the operator aggregates them into a quorum certificate $\m{precommitQC}$ and broadcasts a \txt{PreCommit} message containing it to the ICS. 
% Validators in the ICS commit the tesseract to their history on receiving the message and set their $\m{lockedQC}^p$ as $\m{precommitQC}$ for each interacting participant $p$ whose context they belong to.
% The operator then asynchronously broadcasts the tesseract to all validators, who commit it upon reception. 
% This message includes the quorum certificate, randomly sampled stochastic set, corresponding pseudo-random number, and its proof for validators to verify the message.


% \subsubsection{End of the view}

% At the end of the view, all validators stop participating in consensus for interactions.
% \rishal{Should I remove the following text? I have already mentioned this in sec 5.1.2.}
% Since consensus duration is always a multiple of $\m{viewLength}$, validators enter into achieving consensus on new interactions in sync, ensuring their availability for participating in consensus on future interactions.
% If consensus is not achieved on the proposed tesseract within the allocated time, any future interaction with the same participant set will have its consensus timeout duration increased by one $\m{viewLength}$. This extended duration helps validators combat potential network delays among the participants' contexts.
%This is to make progress during the epoch in the face of network delays among the validators in the interacting participants' contexts. 

% \subsubsection{Conflicting Participant Sets}

% A key difference between Krama and most blockchain consensus protocols is that Krama reaches consensus on individual tesseracts by forming a committee of validators from the participants' contexts.% for each participant and consensus is achieved by forming a set of validators consisting of the interacting participants' contexts. 
% This difference gives rise to the following scenario.
% \begin{example}[Overlapping Participant Sets]\label{ex:conflictingixns}
%     Consider two tesseracts $T_1$ and $T_2$ with participant sets $P_1 = \{p_1, p_2\}$ and $P_2 = \{p_2, p_3\}$.
%     Let the current view be $v$ where $\msc{operator}(P_1, v)\in \calC_{p_1}$ and $\msc{operator}(P_2, v) \in \calC(p_3)$.
%     Here, since the operators for both tesseracts are different, both can propose the respective tesseracts. 
%     However, honest validators in $p_2$'s context can only vote for either of the two tesseracts as they have a common participant. 
%     Hence, it can be the case that neither of the operators obtain a supermajority of prevotes from $p_2$'s context, hence failing to reach consensus on both the tesseracts.
%     Note that this can occur in a view after GST with honest operators for both tesseracts.
% \end{example}

% We see that different operators can propose tesseracts with unequal participant sets and a nonempty intersection in the same view.
% Thus, Krama's normal operation described above has the following modifications to eventually resolve such conflicting proposals resulting in failure to reach consensus. 
% First, we introduce some notation. 
% For a validator $v$, let $P_v$ denote the set of participants in whose contexts the validator belongs to. 
% The power set of $P_v$ is written as $\frakP_v$.

% Over consecutive views, context validators of the overlapping participants arrive at a single tesseract among the conflicting ones for each participant. 
% A validator $v$ on obtaining more than one valid proposal for conflicting tesseracts from different operators, store them in a set $\m{conflicts}_P$ called the \emph{conflict set} for the set of participants $P$.
% Here, $P\in \frakP_v$ is the intersection of the tesseracts' participant sets along with $P_v$. 
% More precisely, $P=\bigcap\limits_{T\in \bar{T}}P_T \cap P_v$ where $\bar{T}$ denotes the tesseracts obtained by the validator in proposals. 
% Additionally, the validator stores the \txt{Propose} message obtained from the operator as justification.
% Note that two conflicting proposals from the same operator signifies Byzantine behaviour by the validator and hence validators ignore the second proposal and take it as proof of Byzantine behaviour for penalisation.
% Validators deterministically obtain a single tesseract from a conflict set by the function $\msc{Tesseract}(\m{conflicts}_P)$.
% %A single tesseract from a participant's conflict set is deterministically obtained by validators by the function $\msc{Tesseract}(\m{conflicts}_p)$. 
% Thus, validators with the same conflict set arrive at the same tesseract. 
% During the Prepare stage for an interaction $I$, in addition to sharing the locked QCs, the validator shares its conflict set for the set of participants $P_I\cap P_v$. 
% \rishal{This is currently being updated.}
% The operator, on receiving them from the interaction's context, combines them into a single conflict set for each such set of participants. 
% The operator shares this as part of the \txt{Propose} message and validators on receiving it, verify it and update their corresponding conflict set for the set of participants $P_I\cap P_v$.
% %Validators verify and update their conflict sets after receiving the proposed message.
% Operators can only Propose the determined tesseract from the set of conflicts. 
% Validators only vote in Prevote and Precommit stages to the tesseract obtained from the set of conflicting ones. 
% Note that all tesseracts in a participant's conflict set stored by a validator have the same height for the participant. 
% When a tesseract for this height for the participant is agreed upon, the validator empties the participant's conflict set. 



% \subsection{Utilities}\label{subsec:utility}
% In this section, we describe the procedures used in Krama to select the stochastic set, and initialise new participants.
% The procedure for stochastic set selection by the operator is given in \Cref{alg:sample}.


% \subsubsection{Formation of the Stochastic Set}\label{subsubsec:randsetsample}







    \subsection{Blockchain Semantics}
\label{sec:blockchain} 

The protocol described above is agnostic to the underlying blockchain semantics. Here we describe the semantics of \emph{interactions} and \emph{Tesseracts} in the context of a blockchain.

\subsubsection{Genesis tesseract}

The genesis tesseract introduces a special participant $\msc{Sarga}$ (Sanskrit for `genesis') that is responsible for maintaining the state of the blockchain. The context of $\msc{Sarga}$ is involved in initialising the context of new participants joining the system. %contains the contexts of all the participants in the system.

\subsubsection{New participants}

%When a new participant $p$ joins the system, it is added to the context of $\msc{Sarga}$ via an interaction that updates the state of $\msc{Sarga}$ to include the context of $p$. When the interaction is committed, the new participant is allowed to be part of further interactions.
When a new participant $p$ joins the system, its context is initialised via an interaction agreed upon by the context of $\msc{Sarga}$.% that updates the state of $\msc{Sarga}$ to include the context of $p$. 
When the interaction is committed, the new participant is allowed to be part of further interactions.

\subsubsection{Context updates between epochs}

\codename{} pre-determines the context of each participant at the start of an epoch. However, when validators fail and no longer participant in \emph{consensus instances}, 
the interactions of the corresponding participants may get blocked due to the lack of a quorum. To mitigate such scenarios, \codename{} allows for special interactions - \emph{context updates} - on $\msc{Sarga}$ to update the context of a participant $p$ in the middle of an epoch. This is done by creating a new interaction that updates the context of $p$ in $\msc{Sarga}$'s state.

Validators when voting on \emph{context updates}, additionally verify that marked validators are indeed faulty or unresponsive.

\subsubsection{Quasi-permissionless}

The reconfiguration protocol allows for validators to be added or removed from the system during epoch transitions, with the entire set of validators $\bbV$ in an epoch considered active. Therefore, \codename{} is characterized as a \emph{quasi-permissionless} blockchain as per the definition in~\cite{LR24}.
    \section{Analysis}
\label{sec:analysis}

In this section, we describe and prove the safety and liveness properties of \codename{}.

To recall, \codename{} accepts interactions $I$ (analogous to transactions) that result in Tesseracts $T$ (analogous to blocks) being committed by validators. We follow the traditional definitions of safety and liveness in the context of blockchains,

To show that \codename{} satisfies both Safety and Liveness, we rely on a key guarantee of the context shuffling and stochastic set sampling mechanisms used in the protocol.

\begin{assumption}[Random Sampling Assumption]
    The stochastic set sampling and context shuffling mechanisms used in \codename{} ensure that with high probability, the sampled set of validators contains strictly less than one-third Byzantine validators.
\end{assumption}

We satisfy the above assumption by choosing appropriate sizes for the contexts, stochastic sets, and the validator pool and demonstrate that it is met with high probability by modelling the sampling procedure using a hypergeometric distribution. We refer the reader to the Appendix for a detailed probabilistic analysis.
%We satisfy the above assumption by using a hypergeometric distribution to sample a fixed-size set of validators from the total validator pool. We refer the reader to the Appendix for a detailed probabilistic analysis.

% \subsection{Guarantees}
% We now define the properties achieved by Krama with respect to the histories stored by validators. 

% The following properties are achieved by Krama when all participant contexts and sampled stochastic sets have less than one-third Byzantine validators. 
% This setting is obtained with high probability as demonstrated below in \Cref{subsec:probanalysis}.

% \paragraph*{Safety} 
% Two conflicting tesseracts cannot both be committed by honest validators.

% \paragraph*{Liveness}
% In each view beginning after GST, in each maximal set of overlapping consensus instances with honest operators, at least one interaction is successfully committed by all honest validators within a bounded time $T$.
% % \\ In each view after GST, in every set of overlapping interactions for which Prepare messages are sent, at least one interaction is successfully committed by all honest validators.  
% % \\ After GST, in every set of consensus instances with overlaps where each of the instances is led by an honest operator, at least one of them successfully leads to a consensus decision.
% % \\ After GST, every proposal by honest operators or a conflicting tesseract is eventually committed by all honest validators. 

% The rest of this section is devoted to showing that the above guarantees are achieved. 

% \subsection{Random Sampling}\label{subsec:probanalysis}
% We first provide a probabilistic analysis of random sampling used in Krama. 
% We randomly sample a subset of validators in two key areas: (1) stochastic set formation, and (2) context updation.
% In both instances, the fundamental process involves randomly sampling a fixed-size set of validators without replacement from the total validator pool.
% We demonstrate that, with high probability, the sampled set of validators contains strictly less than one-third Byzantine validators.

% Consider a set of $N$ validators with $B$ Byzantine faults. 
% The number of Byzantine validators $b$ in a random sample of $n$ validators from the above can be given as follows.
% The probability distribution of the number of Byzantine faults in the sampled set is given by the hypergeometric distribution. The hypergeometric distribution is appropriate in this context as it models the probability of selecting a certain number of "successes" (in this case, Byzantine validators) from a finite population without replacement. 
% The probability mass function of the hypergeometric distribution for this analysis is given by:
% $$P(b = k) = \frac{ {B \choose k} {N-B \choose n-k} }{ {N \choose n} }$$
% The probability that the stochastic set has less than one-third Byzantine faults can be calculated by using the cumulative distribution function of the hypergeometric distribution. On calculating for a network of 1000 validators, of which one-fifth are Byzantine (according to our assumptions), for $n=67$, the required probability is $0.99682$. 
%The context power used as a weight for each validator in the network can be assumed to be reasonably uniformly distributed among the validators. This is ensured by tokenomics, restricting accumulation of context power by a single individual. Thus, it follows that with high probability, the stochastic set chosen has less than one-third of its context power held by Byzantine validators. 

% Thus, the assumptions made in the protocol to ensure the protocol satisfies safety and liveness are met with high probability for the above parameters.
% In our implementation, the size of participants' contexts and the stochastic set, $n_\calC = 67$ and hence the above analysis accurately captures the protocol's implementation parameters.
% In the following analysis, we assume that this is the setting in which the protocol is operating.

\subsection{Safety}
%The analysis of safety of \codename{} is inspired by the analysis of Fast-Hotstuff's safety in \cite{JNFG23}.
This section is devoted to proving the following theorem.

\begin{theorem}[Safety]\label{thm:safety}
    Between two conflicting Tesseracts $T$ and $T'$, every honest validator will commit at most one of them.
\end{theorem}

Starting from the quorum intersection property, we first show that only one of two conflicting Tesseracts are voted for by honest validators. Then, we show that once a \emph{supermajority} of validators vote for a Tesseract, it is guaranteed that the Tesseract will be committed in either this or later rounds therefore preserving safety.

% In this section, we prove safety of the Krama protocol. 
% The structure of our analysis is inspired by Fast-Hotstuff's analysis given in \cite{JNFG23}.

\begin{lemma}[Quorum Intersection Property]\label{lem:qcintersection}
    Let $S$ be a set of $n=3f+1$ validators, with at most $f$ Byzantine faults.
    For any two supermajorities over the set of validators $S$, we have at least one honest validator in common. 
\end{lemma}


% \begin{proof}
%     Let $M_1$ and $M_2$ be two supermajorities over $S$. 
%     Since a vote supermajority consists of $n-f$ validators by definition, $2(n-f)=2n-2f=n+f+1$. Therefore, any two supermajorities have at least $f+1$ validators in common. As the number of Byzantine faults is at most $f$, the result follows. 
% \end{proof}

\begin{lemma}\label{lem:conflictdiffview}
    If there are valid $\m{voteQCs}$ for conflicting tesseracts $T$ and $T'$ then they were formed during different views.
\end{lemma}
\begin{proof}
    We prove by contradiction. 
    Suppose, there exist two QCs $\m{voteQC}$ and $\m{voteQC}'$ for the tesseracts $T$ and $T'$ formed during the same view.
    Since the tesseracts $T$ and $T'$ are conflicting, there exists a participant $p\in P_T\cap P_{T'}$.
    Therefore, by definition of a Tesseract's ICS, $\calC_p\subseteq ICS(T)$ as well as $\calC_p\subseteq ICS(T')$. 
    Now, by \Cref{lem:qcintersection}, we have an honest validator $v\in \calC_p$ that is a part of both QCs.
    However, honest validators do not vote for multiple tesseracts involving the same participant in a given view. 
    Hence, the result follows by contradiction.
\end{proof}


\begin{lemma}\label{lem:highestqc} 
    If an honest validator has received $\m{commitQC}$ for tesseract $T$ then the corresponding $\m{voteQC}$ will be a $lockedQC$ part of the \txt{Prepare} messages obtained by the operator in the next consensus instance for each participant $p\in P_T$.
\end{lemma} 
\begin{proof}
    If an honest validator received a $\m{commitQC}$, then the validators whose \txt{PreCommit} messages comprise the quorum certificate have received $\m{voteQCs}$. % was received by an honest validator in the view, a quorum of validators from its ICS had sent \txt{PreCommit} messages. 
    Therefore, a quorum of validators are locked onto the $\m{voteQC}$. 
    Now, consider the next view with a consensus instance containing a participant $p\in P_T$.
    During the Prepare stage, say the operator obtains a supermajority of \txt{Prepare} messages from $p$'s context.
    Note that the quorum of validators locked onto the $\m{voteQC}$ consists of a supermajority of votes from $p$'s context.
    Therefore, by \Cref{lem:qcintersection} there exists an honest validator $v\in\calC_p$ that is present in both of the above supermajorities.
    %Now, for the next view, the operator will receive a quorum of \txt{Prepare} messages containing $\m{lockedQC}$s in the prepare stage. 
    %Since a quorum contains a supermajority of votes from each interacting participant $p$'s context, by \Cref{lem:qcintersection} there exists an honest validator $v\in\calC_p$ that is present in both of the above quorums.
    This validator sends the $\m{voteQC}$ as its locked QC in the \txt{Prepare} message and is therefore obtained by the operator. 
\end{proof}
Thus, tesseracts previously committed by honest validators are learned by operators during the prepare stage at the start of each view.

\begin{lemma}\label{lem:safetyinduction}
%If an honest validator has committed a tesseract $T$ in view $v$, then every tesseract $T'$ such that $P_T\cap P_{T'}\neq\emptyset$, for which a $voteQC$ is obtained in view $v'>v$ is a descendant of $T$. 
Let $T$ be a tesseract committed by an honest validator in view $v$. For every tesseract $T'$ with a $\m{voteQC}$ formed in view $v'>v$ where $P_T\cap P_{T'}\neq\emptyset$, we have that $T'$ is a descendant of $T$, i.e., $T\preceq T'$.
\end{lemma}
\begin{proof}
We prove by induction on the views $v'>v$. 

\noindent{Base Case:} 
Let $v'=v+1$. Since an honest validator has committed tesseract $T$, by \Cref{lem:highestqc}, the corresponding $\m{voteQC}$ is obtained as part of the $\m{prepareQC}$ during the \txt{Prepare} stage in the consensus instance of tesseract $T'$. 
There are two outcomes for this $\m{voteQC}$: either it is the unique \txt{Vote} stage QC present in $\m{prepareQC}$ and thus reproposed, or it is removed by the \textsc{RemoveOutdated}() procedure. 
In the former case, the tesseract $T'$ is simply a reproposal of the tesseract $T$ and is thus is a descendant. 
Otherwise, if the $\m{voteQC}$ is removed, then there exists a higher $\m{lockedQC}$ obtained in the \txt{Prepare} stage. However, this QC must be from a view $v_l<v$ as there can be no other $\m{voteQC}$ other than that of $T$ from view $v$ due to \Cref{lem:conflictdiffview}. 
If this $\m{lockedQC}$ is from \txt{Commit} stage in the same height as $T$ for a common participant or is from \txt{Vote} stage with a higher height, then it both cases, a validator from the common participant's context has committed a tesseract in the same height as $T$. 
Then, this tesseract would have been discovered in future views until view $v$. 
This tesseract thus has a quorum of validators to be locked onto it. This tesseract therefore has 
It must be the case that this tesseract is $T$ due to \Cref{lem:conflictdiffview}. 


If this $\m{lockedQC}$ is from \txt{Vote} stage 
then it is not higher than the $\m{voteQC}$ of tesseract $T$ as it is from a later view. 
Therefore, it has to be a \txt{Commit} stage QC. 
However, 
In the \textsc{RemoveOutdated}() procedure, this $\m{voteQC}$ 

is not removed due to comparison with a higher tesseract as the views are in succession and hence such a tesseract could not have been formed. Let, if possible, a higher tesseract was formed 

\noindent{Inductive Case:}
Now, we assume that the lemma holds upto view $v'-1>v$ and prove that it also holds for view $v'$. 
By the induction hypothesis, $\m{lockedQC}$s obtained in the \txt{Prepare} stage from participants in $P_T\cap P_{T'}$ correspond to tesseracts that are descendants of $T$. 
For each such $\m{lockedQC}$ from the \txt{Vote} stage, there are two outcomes: it can be reproposed as it is the unique \txt{Vote} stage QC present (thus corresponding to $T'$), or it can be removed by the \textsc{RemoveOutdated}() procedure. 
In the former case, the lemma holds by induction hypothesis.
If the QC is removed, then there exists a tesseract higher than the one corresponding to the $\m{lockedQC}$, say $T_l$. 
Such a tesseract is also a descendant of $T$ by induction hypothesis as must be formed in a view $v_l$ between $v$ and $v'-1$, as otherwise it would not be higher than $T$ and thus not higher than $T_l$. 
Now, if this tesseract $T_l$ is also removed due to the presence of a higher tesseract, then that tesseract also shall satisfy the induction hypothesis. Therefore, any tesseract that is reproposed or extended shall be a descendant of $T$ by induction hypothesis and hence the result follows.
\end{proof}

Now, from \Cref{lem:safetyinduction} and  \Cref{lem:conflictdiffview} we obtain \Cref{thm:safety}. 

%%%Old Proof - Incorrect%%%
%\begin{theorem}
%    Two conflicting tesseracts cannot both be committed by honest validators.
%\end{theorem}
%\begin{proof}
%    An honest validator commits a tesseract only upon receiving a valid $\m{commitQC}$. 
%    From \Cref{lem:conflictimpliesdiffslot}, it follows that the conflicting tesseracts' $\m{commitQC}$s were formed in different views. 
%    Thus, without loss of generality, consider two conflicting tesseracts $T$ and $T'$ with a participant $p\in P_T\cap P_{T'}$ where $\msc{view}(T)<\msc{view}(T')$. 
%    For tesseract $T'$ to conflict with $T$, it has to extend an ancestor of $T$. 
%    However, by \Cref{lem:highestqc}, we know that the operator for the subsequent view with an interaction for participant $p$ after $\msc{view}(T)$ will have the $\m{voteQC}$ formed for tesseract $T$ as part of its $\m{prepareQC}$
%    The operator, while determining which interaction to propose, thus considers the above $\m{voteQC}$ as well. 
%    In this procedure by the operator, there are two outcomes for this QC: it can be removed from the set of locked QCs or it can stay as part of the set. %, the operator can either remove this QC be removed from the set of locked QCs or not. 
%    If it is not removed, then the proposal has to be a $p$-child to tesseract $T$ based on the protocol rules that proposals must satisfy. 
%    If it is removed, then the proposal has to extend the latest state of each of the interacting participants, and therefore it will be a $p$-child of the committed tesseract $T$.
%    Thus in both cases the resulting proposal is a $p$-child of $T$ and therefore $T'$ cannot extend an ancestor of $T$.
%    Moreover, operators cannot forge a $\m{prepareQC}$ for $T'$ since validators verify the quorum certificates. 
%    Hence, the result follows.
%\end{proof}

\subsection{Liveness}
\label{subsec:liveness}

In this section, we prove liveness of \codename{}.

A key novelty of \codename{} is the use of consensus instances that are updating disconnected parts of the global states - or participant sets. Therefore, a global liveness property should ensure first that among the many overlapping consensus instances, at least one succeeds in committing a Tesseract.

To illustrate, consider a set of overlapping consensus instances, each with a different set of participants. As described in~\cref{fig:overlaps}, we have 4 concurrent consensus instances with participants $p,q,r,$ and $s$ -- $\{p,q\}, \{q,r\}, \{p,r\},$ and $\{r,s\}$. Here, the consensus instances $\{p,q\}$, $\{q,r\}$, and $\{p,r\}$ overlap with each other. Yet, \codename{} cannot ensure that among these set of overlapping consensus instances, at least one instance will successfully commit. Since the validators associated with $p$, $q$, and $r$ can choose to respond to the consensus instance $\{p,q\}$, $\{q,r\}$, or $\{p,r\}$, they may not all agree on a single instance required to achieve a quorum.

However, the maximal set of overlapping consensus instances is $\{p,q\}, \{q,r\}, \{p,r\},$ and $\{r,s\}$, as the instance $\{r,s\}$ overlaps with $\{q,r\}$ and $\{p,r\}$. In this case, \codename{} ensures that at least one of the consensus instances in the maximal set will successfully commit a Tesseract. Recall that in the \txt{Prepare} phase of \codename{}, validators pick a prepare message to respond from the set of overlapping consensus instances using the $\msc{Choose}()$ function. By ensuring that the $\msc{Choose}()$ function deterministically orders the consensus instances and picks the highest one among them, we can guarantee that at least one consensus instance in a maximal set of overlapping consensus instances will be chosen by all participants in the context of the participants in the maximal set.

Formally, we define overlapping consensus instances and maximal sets of overlapping consensus instances as follows.

\begin{definition}[Overlapping Consensus Instance]
    A consensus instance is said to be overlapping with another if they share at least one participant in their participant sets.
\end{definition}

\begin{definition}[Maximal Set of Overlapping Consensus Instances]
    A maximal set of overlapping consensus instances is a set of consensus instances such that every instance in the set overlaps with at least one other instance in the set, and no additional instance can be added to the set without violating the overlap condition.
\end{definition}

% \begin{example}\label{ex:overlaps}
% 	Consider a set of four participants $p,q,r$ and $s$ with consensus instances defined by the participant sets $\{p,q\}, \{q,r\}, \{p,r\},$ and $\{r,s\}$ (See \Cref{fig:overlaps}).
% 	Here, for instance, the consensus instances given by $\{p,q\}$ and $\{q,r\}$ overlap as they have $q$ in common. 
% 	The set of consensus instances given by $\{p,q\}, \{q,r\},$ and $\{p,r\}$ is not maximal here as the consensus instance with $\{r,s\}$ can be added as it overlaps with $\{q,r\}$ and $\{p,r\}$.
% \end{example}

\begin{figure}
\tikzfig{overlaps}
\caption{A maximal set of overlapping consensus instances. Vertices represent participants and edges represent consensus instances. Here, the set of black consensus instances is not maximal.}
\label{fig:overlaps}
\end{figure}

% Toward guaranteeing that at least a single consensus instance achieves a quorum among a maximal set of overlaps, the $\msc{Choose}()$ function used by validators is as follows.
% Among a set of overlapping consensus instances, $\msc{Choose}()$ deterministically orders the consensus instances and picks the highest one.
% To demonstrate how this achieves the guarantee, consider a maximal set of overlapping consensus instances. 
% Here, each participant's context, chooses from the subset of instances containing the participant. 
% Therefore, among a maximal set of overlaps, the highest instance in the order defined by $\msc{Choose}()$ is chosen by each of its participants' contexts.  
% However, this is not true among a non-maximal set of overlaps.
% For instance, in \Cref{ex:overlaps}, let the order among the consensus instances be $I_{pq}< I_{qr}< I_{pr}<I_{rs}$.
% Here, validators in the context of participants $p, q, r$, and $s$ choose $I_{pr}, I_{qr}, I_{rs},$ and $I_{rs}$ respectively. 
% Note that among the non-maximal set of instances given in black, none are chosen by both its participants.
% This leads to the below result. 

The rest of the section is devoted to proving the following theorem.

\begin{theorem}[Liveness] 
    In each view beginning after GST, in each maximal set of overlapping consensus instances with honest operators, at least one interaction is successfully committed by all honest validators within a bounded time $\tau$.
\end{theorem}



Prior to proving liveness, we first show that in each maximal set of overlapping consensus instances with honest operators, at least one consensus instance successfully obtains a quorum of \txt{Prepare} messages from each of its participants' contexts. Subsequently, we use the lemma to show that at least one interaction among a maximal set will eventually be successfully committed by all honest validators. 

\begin{lemma}\label{lem:prepareOverlap}
    In each maximal set of overlapping consensus instances with honest operators, at least one consensus instance successfully obtains a quorum of \txt{Prepare} messages from each of its participants' contexts. 
\end{lemma}
\begin{proof}
    Consider a view after GST with a maximal set of overlapping consensus instances led by honest operators. 
    Let $P$ be the union of all the participant sets in the consensus instances. 
    Validators in each context $\calC_p$ for $p\in P$, choose a single consensus instance from the \txt{Prepare} messages they have received. 
    Let $\m{prepares}_p$ be the set of \txt{Prepare} messages received by validators in $\calC_p$.
    Since the operators are honest, all validators in a participant's context receive the same set of \txt{Prepare} messages as operators multicast them to the entire context. 
    Validators make the choice of a consensus instance to respond to by ordering the \txt{Prepare} messages that they have received. 
    This order is based on view randomness and is thus common to all context validators for participants in $P$. 
    Consider the consensus instance that is the maximum element in the order. 
    Each of its participant set context validators will choose it as it is the highest among the \txt{Prepare} messages they have received as it is the maximum element. 
    Further, since the set of overlapping consensus instances is maximal, there is no other consensus instance that intersects with any of the instances in the given set. 
    Therefore, the consensus instance that is maximum in the order successfully obtains a quorum.
\end{proof}

\begin{algorithm}[hptb]
\caption{Removing outdated locks}\label{alg:outdated}
\begin{algorithmic}[1]         

\Procedure{RemoveOutdatedQCs}{$\m{qcSet}$}
     \State $\msc{ComputeHighQCs}(\m{qcSet})$
     %\State $\{\m{maxHeightPrecommit}_p\mid \exists \m{qc}\in \m{qcSet} : p\in \m{qc}.\m{participants}\} \gets \msc{MaxHeight}(\m{qcSet}_v, \m{precommit})$
     %\State $\m{maxHeightPrevote}_v \gets \msc{MaxHeight}(\m{qcSet}_v, \m{prevote})$
     \ForAll{$\m{qc} \in \m{qcSet}$ such that $\m{qc}.\m{stage} = \m{vote}$}
         \ForAll{$p \in \m{qc}.\m{participants}$}
             \If{$\m{qc}.\m{height}(p) \leq \m{highCommitQC}_p.\m{height}(p) \lor \m{qc}.\m{height}(p) < \m{highVoteQC}_p.\m{height}(p)$}
                 \State $\m{qcSet} \gets \m{qcSet} \setminus \{\m{qc}\}$
                 \State \textbf{break}
             \ElsIf{$\m{qc}.\m{height}(p) = \m{highVoteQC}_p.\m{height}(p) \land \m{qc}.\m{view} < \m{highVoteQC}_p.\m{view}$}
                 \State $\m{qcSet} \gets \m{qcSet} \setminus \{\m{qc}\}$
                 \State \textbf{break}
             \EndIf
         \EndFor
     \EndFor
 \EndProcedure
\item[]

\end{algorithmic}
\end{algorithm}


Additionally, during the \txt{Prepare} stage for a given participant set $P$, Byzantine validators can send outdated \txt{Vote} locks in whose place other tesseracts have been committed with participants outside $P$ to prevent the operator from making a proposal due to the presence of more than one \txt{Vote} locked tesseract. 
Therefore, it is crucial for liveness that such outdated information is removed from the locks obtained. 
This is accomplished by the $\msc{RemoveOutdated}()$ procedure given in \Cref{alg:outdated} that we describe below.
The outdated locks are removed by comparing with the maximum height locks obtained for each participant present in the locked tesseracts. 
The outdated vote QCs are filtered by removing those where:
     \begin{enumerate}
         \item A participant $p$ in the QC has a commit QC with atleast the same height for the participant
         %the same height or higher in some precommit stage QC by comparing with $maxHeightPrecommit_p$ for the participant, or
         \item A participant $p$ in the QC has a vote QC higher for the participant
         %higher height in some other prevote stage QC by comparing with $maxHeightPrevote_p$.
     \end{enumerate}
This ensures that vote locks that have been superseded by committed tesseracts are ignored.

\begin{theorem} 
    %In a view after GST, every tesseract proposed by honest operators, or another conflicting tesseract is eventually committed by all honest validators. 
    In each view beginning after GST, in each maximal set of overlapping consensus instances with honest operators, at least one interaction is successfully committed by all honest validators within a bounded time $T$.
\end{theorem}
\begin{proof}
    By \Cref{lem:prepareOverlap}, in the given scenario, at least one consensus instance obtains a quorum of \txt{Prepare} messages and therefore proceeds to make a proposal. 
    The outdated locked QCs received in the quorum of \txt{Prepare} messages are removed by the operator using the $\msc{RemoveOutdatedQCs}()$ procedure.
    This ensures that Byzantine validators cannot prevent proposals being made by sending outdated locks for tesseracts with different participant sets.
    Therefore, a tesseract extending the latest locks is proposed as per the $\msc{ExtendCommitted}()$ function ensuring safety. 
    During this view, validators in the ICS can successfully commit the proposal as supermajorities of \txt{Vote} and \txt{Commit} messages are formed during the respective stages, since message delay is bounded by $\Delta$ and the length of the view is assumed to be sufficiently long for the voting rounds' communication to occur between the operator and validators. 
    Here, supermajority of votes are formed as validators send \txt{Vote} messages if the proposed tesseract is the result of the proposal determining procedure run by the operator, as given in \Cref{validpropose}
        Validators not in the ICS commit the tesseract upon receiving it from the operator.
\end{proof}

    \section{Related Work}
\label{sec:relatedwork}

% Github copilot generated this section, with some tweaking, from the references.

Recent years have witnessed significant advancements in distributed ledger technologies, particularly in the design of consensus protocols and scalable blockchain architectures. Early systems such as Bitcoin~\cite{Nakamoto09} introduced decentralized payment networks, while Ethereum~\cite{But15} extended this vision to general-purpose smart contracts, enabling a broader range of decentralized applications. As the adoption of blockchains has grown, so has the need for scalable, secure, and decentralized consensus mechanisms.

Consensus protocols in blockchains can be broadly categorized into longest-chain protocols, such as those in the Ouroboros family~\cite{KRDO17, DGKR18, BGKRZ18} and Snow White~\cite{DPS19}, and protocols inspired by classical Byzantine Fault Tolerance (BFT) research~\cite{Bra87, LSP82, CL99}. The latter category includes modern protocols like Tendermint~\cite{BKM18}, Algorand~\cite{GHMVZ17}, Hotstuff~\cite{YMRGA19}, and Fast-Hotstuff~\cite{JNFG23}, primarily motivated by PBFT~\cite{CL99}, which leverage round-based consensus with rotating leaders to improve performance and robustness. Avalanche~\cite{Roc20} represents a distinct approach, employing metastable consensus through repeated sampling without a designated leader. Comprehensive surveys such as~\cite{BS+19} provide an overview of these developments.

AptosBFT~\cite{Aptos22} and Sui~\cite{Sui22} are recent protocols that extend the Hotstuff family, focusing on high throughput and low-latency consensus by leveraging parallel execution and advanced mempool designs. Both systems employ a global consensus protocol to order transactions and achieve scalability by decoupling execution from consensus, allowing for efficient processing of independent transactions. In contrast, \codename{} departs from the global ordering paradigm by associating a set of validators with each participant and maintaining a separate chain per participant. Rather than relying on a single global consensus instance, \codename{} commits multi-participant transactions across the relevant chains, enabling parallelism without requiring a total order on all transactions.

Committee-based consensus protocols have emerged as a means to enhance scalability and efficiency. In these systems, a subset of validators (the committee) is selected to participate in consensus for a given period or instance~\cite{ZMR18, GHMVZ17, PS17}. Dynamic committee formation, as adopted in our protocol, allows multiple consensus instances to proceed in parallel, further improving throughput.

Sharding is another prominent technique for scaling blockchains, wherein the network is partitioned into multiple shards, each processing a subset of transactions independently~\cite{KJ+17, SPW19, Zil17, Har}. While sharding increases throughput, it introduces challenges in cross-shard communication and state management. In contrast, our approach associates a set of validators with each participant, maintaining a separate chain per participant and committing multi-participant transactions across relevant chains. Unlike state sharding, all validators in our protocol maintain the entire system state, simplifying consistency and security guarantees.

Finally, while Krama employs a multi-linked Directed Acyclic Graph (DAG) structure to record validator states, it differs fundamentally from DAG-based consensus protocols such as DAG-Rider~\cite{KKNS21}, Bullshark~\cite{SGSK22}, and Cordial Miners~\cite{KNPS22}, which decouple block dissemination from total ordering. In Krama, the DAG serves as a local log of states to be agreed upon, rather than as a mechanism for achieving consensus itself.

% Additional related works:
% - HoneyBadgerBFT~\cite{Miller16} for asynchronous BFT consensus.
% - Red Belly Blockchain~\cite{RedBelly18} for scalable BFT consensus.
% - Chainspace~\cite{Al-Bassam18} for sharded smart contract platforms.
% - Elastico~\cite{Luu16} as an early sharding protocol.
% - Spanner~\cite{Corbett12} and Google’s TrueTime for global consistency in distributed systems.
% - Hyperledger Fabric~\cite{Androulaki18} as a permissioned blockchain with pluggable consensus.
% - Libra/DiemBFT~\cite{Baudet20} for production-grade BFT consensus in payment systems.
% - Narwhal and Tusk~\cite{Danezis22} for separating mempool and consensus layers.
% - Sui~\cite{Sui22} and Aptos~\cite{Aptos22} for recent advances in parallel execution and consensus.

% \subsection{Blockchains and Consensus}
% - Vision started with payment systems - Bitcoin \cite{Nakamoto09}. 

% - Then a general purpose smart contracts platform - Say Ethereum whitepaper from 2014/15 \cite{But15}.

% - Adding these above because part of the vision behind Krama is that day to day interactions will be done more in the blockchain and hence more validators joining provides more decentralised trust (Perhaps refer to \cite{DMR20} - not sure how relevant of a paper this is - unfamiliar with that side of literature) 

% - A good survey paper from 2019 is \cite{BS+19}. \rishal{not aware of a similar paper from more recent times.}

% - BFT consensus protocols for blockchains are of largely two types 

% (1) longest chain protocols inspired by nakamoto consensus - Ouroboros family \cite{KRDO17, DGKR18, BGKRZ18}, Snow White \cite{DPS19} 

% (2) protocols inspired by BFT research from the 20th century for a long time with works such as \cite{Bra87}, \cite{LSP82}, PBFT \cite{CL99} - examples include Tendermint\cite{BKM18}, Algorand\cite{GHMVZ17}, Hotstuff\cite{YMRGA19}, Fast-Hotstuff\cite{JNFG23} - typically round based with a rotating leader 

% One interesting family of consensus protocols here is Avalanche\cite{Roc20} which is leaderless and uses repeated sampling to reach a decision - 'metastable consensus'.

% - Our protocol is also in this line of 'BFT-type' protocols, related to the Hotstuff family of protocols, specifically Fast-Hotstuff.

% - Write about Fast-Hotstuff and its two phase normal non-pipelined protocol and how it uses proof of highest QC in the proposal to improve over Hotstuff's three phases during normal operation.

% Among such protocols, in cases, a committee is randomly sampled and the decision is reached among this global committee 

% \subsection{Committees}

% - Committee based consensus - A good survey of this is done in the literature overview of \cite{ZMR18}. Other examples include Algorand\cite{GHMVZ17}, Hybrid Consensus \cite{PS17}

% Our protocol uses committees formed dynamically for each consensus instance several of which can run in parallel. 

% \subsection{Shards}

% One approach to increase throughput as validators increase is to have multiple fixed committees of nodes called shards that process transactions in parallel in different chains with each shard having its own chain whose txns it processes - throughput increase being a multiple (\# of shards) of what it originally was per shard. 

% - Sharding - Examples include Omniledger\cite{KJ+17},  Nightshade\cite{SPW19} - Near protocol's sharding, Zilliqa\cite{Zil17}, Harmony\cite{Har} whitepapers
% Cross shard communication etc gets really complicated - No issues like that in Krama. 
% Can only shard communication and transaction validation or shard the blockchain state as well in addition. 
% Participant agnostic. 

% In Krama, we have a set of validators associated with each participant, with each participant having its own chain and tesseracts with multiple participants being committed in each of their chains unlike sharding where transactions are only processed in its relevant shard. 
% Krama's consensus naturally involves each of these sets associated with participants. 
% Further, all validators maintain the entire state unlike state sharding. 


% \subsection{DAG-based Consensus}
% The multi-linked DAG in Krama is not to be confused with DAG based consensus protocols such as DAG-Rider, Bullshark\cite{SGSK22}, and Cordial Miners\cite{KNPS22} that separate the block dissemination mechanism from the task of totally ordering the chain. (\rishal{unsure of latest work in DAG based consensus. I am not very familiar with this in general. }) 
% The multi-linked DAG is the log of states stored by each validator that is to be agreed upon by the protocol. 


    \input{sec7-conclusion}
    \bibliography{references}
\end{document}
\endinput
